\documentclass{book}
%\documentclass{book}
%If you're going to print this, perhaps get ride of oneside so that \parts{} will have their own page!
\def\title{Ordinary Differential Equations}
\def\author{Nathanael Chwojko-Srawley}
%\date{} % If you want to specify a date
\usepackage{inputenc}         % Used to compile any UTF-8 Character (and supports Unicode)
\usepackage{extarrows}
\usepackage{multicol} % For multiple columns
\usepackage{stackengine}
\usepackage{scalerel}
\usepackage{stackrel}         % for left-right arrows, staking symbol above and bellow
\usepackage{mathrsfs}
\usepackage{euscript}
\usepackage{makeidx}
\usepackage[font=itshape]{quoting}
\usepackage{mathtools}
\usepackage{enumitem}
\usepackage{graphicx}         % Let's you use images in LaTeX; ex the \includegraphics command.
\usepackage{wasysym}          % emoji's!
\usepackage{dsfont}
\usepackage{float}            % package with ``H'' for figures and tables, ex. Images, tables, etc
\usepackage{setspace}         % More flexibility on spacing in document.
\usepackage{commutative-diagrams}
\usepackage{hhline}
\usepackage{tikz-cd}          % for commutative diagrams https://tex.stackexchange.com/questions/419221/how-to-do-the-pushout-with-universal-property
\usepackage{longtable}        % create tables that extend past one page
\usepackage{microtype}        % makes nice text. Fixes some under-/over- filled box problems.
\usepackage{xcolor}           % Colour text.
\usepackage{wrapfig}          % To wrap text around figure
\usepackage{listings}         % Create nice colour code blocks (customized in preamble)
\usepackage{amsthm}           % Used to create theorem enviroments.
\usepackage{array}            % \begin{table}{>{\em}c c} -> 1st column italic (bfseries for bold)
\usepackage{blindtext}
\usepackage[a4paper, total={6in, 8in}]{geometry}
\usepackage{fancyhdr}
\usepackage{titlesec}
\usepackage[answerdelayed]{exercise}
\usepackage{etoolbox}          % for Dynkin diagrams
\usepackage{dynkin-diagrams}   % for Dynkin diagrams
\usepackage{pgfplots}

\usepackage{amsmath, amssymb} % Used to write better mathematical expressions (ex. \mathbb{R})
\usepackage[most]{tcolorbox} 	      % Makes nice boxes for thms, cor, prop, or anything else.
\usepackage{thmtools}         % Customize theorem environment (in preamble).
\usepackage{hyperref}
\usepackage{cleveref}


\pgfplotsset{compat=1.18}

% \usepackage{CJKutf8} % for some chinese

% \usepackage{dutchcal} % for lower case mathcal


% For citing
\usepackage{csquotes}
\usepackage{comment}
% \usepackage[backend=biber,citestyle=alphabetic,bibstyle=authortitle]{biblatex}
\usepackage{biblatex}

\addbibresource{algebra_book_cite.bib}

\providecommand{\HUGE}{\Huge}% if not using memoir
\newlength{\drop}% for my convenience
%% select a (FontSite) font by its font family ID
\newcommand*{\FSfont}[1]{\fontencoding{T1}\fontfamily{#1}\selectfont}
%% if you don’t have the FontSite fonts either \renewcommand*{\FSfont}[1]{}
%% or use your own choice of family.
%% select a (TeX Font) font by its font family ID
\newcommand*{\TXfont}[1]{\fontencoding{T1}\fontfamily{#1}\selectfont}
%% Generic publisher’s logo
\newcommand*{\plogo}{\fbox{$\mathcal{PL}$}}

%% Some shades
\definecolor{Dark}{gray}{0.2}
\definecolor{MedDark}{gray}{0.4}
\definecolor{Medium}{gray}{0.6}
\definecolor{Light}{gray}{0.8}
%%%% Additional font series macros
\makeatletter
%%%% light series
%% e.g., kernel doc, section s: line 12 or thereabouts
\DeclareRobustCommand\ltseries
{\not@math@alphabet\ltseries\relax
\fontseries\ltdefault\selectfont}
%% e.g., kernel doc, section t: line 32 or thereabouts
\newcommand{\ltdefault}{l}
%% e.g., kernel doc, section v: line 19 or thereabouts
\DeclareTextFontCommand{\textlt}{\ltseries}
% heavy(bold) series
\DeclareRobustCommand\hbseries
{\not@math@alphabet\hbseries\relax
\fontseries\hbdefault\selectfont}
\newcommand{\hbdefault}{hb}
\DeclareTextFontCommand{\texthb}{\hbseries}
\makeatother


\newcommand{\titleUL}{\begingroup% University of Liege
\drop=0.1\textheight
\vspace*{\drop}
\begin{center}
% {\LARGE\textsc{University of Toronto}}\\[\drop]
% University logo
% {\LARGE \plogo}\\[\drop]
{\LARGE {\quad}}\\[\drop]

\rule{\textwidth}{1pt}\par
\vspace{0.5\baselineskip}
{\huge\bfseries \title\\[0.2\baselineskip]
\large --- \emph{EYNTKA} Series ---}\\[0.5\baselineskip]
\rule{\textwidth}{1pt}\par
\vfill
{\Large\textsc{\author}}
\vfill
\vfill
{\large \today}
\end{center}
\endgroup}

%header information
\pagestyle{fancy}
\setlength\headheight{20pt}
\renewcommand{\sectionmark}[1]{\markright{\thesection.\ #1}}
\renewcommand{\chaptermark}[1]{\markboth{#1}{}}
\fancyfoot[C,CO]{\textcolor{gray}{Nathanael Chwojko-Srawley}}
\fancyfoot[R,RO]{\thepage}{}

\titleformat
{\chapter} % command
[display] % shape
{\bfseries\Huge\itshape} % format
{\thechapter} % label
{0.5ex} % sep
{
    \rule{\textwidth}{1pt}
    \vspace{2ex}
    \centering
} % before-code
[
\vspace{-0.5ex}%
\rule{\textwidth}{1pt}
] % after-code

% Create a new page style for the first page
\fancypagestyle{firstpage}{
  \fancyhf{} % Clear header and footer
  \renewcommand{\headrulewidth}{0pt} % This removes the header line
  \fancyfoot[C]{\thepage}
}

% Apply the first page style conditionally
\AtBeginDocument{\thispagestyle{firstpage}}


\stackMath %to be able to stack text in math mode
% For devnagari font for that one example




%\ExerciseLevelInToc
\counterwithin{Exercise}{section}
\counterwithin{Answer}{section}
%\renewcommand{\ExerciseHeader}{\noindent\def\stackalignment{c}% code from https://tex.stackexchange.com/a/195118/101651
    %\stackunder{{\textbf{\large\ExerciseName}}}{\textcolor{black}{\rule{40pt}{4pt}}}\medskip}



%theorem code
\definecolor{dkgreen}{rgb}{0,0.6,0}
\definecolor{gray}{rgb}{0.5,0.5,0.5}
\definecolor{mauve}{rgb}{0.58,0,0.82}
\definecolor{lightyellow}{rgb}{1,1,0.6}
\definecolor{reddish}{HTML}{A01A3A}
%Theorems, propositions, Corollaries, Definitions

% breakable,

%Wanna implement this, but not working: https://tex.stackexchange.com/questions/547979/problem-numbering-environment
\newtcbtheorem[number within=section]{thm}{Theorem}%
{
colback=green!5,
colframe=green!35!black,
fonttitle=\bfseries,
label type=theorem
}{th}

\newtcbtheorem[number within=section, use counter from=thm]{lem}{Lemma}%
{
colback=blue!5,
colframe=black!35!black,
fonttitle=\bfseries,
label type=lemma
}{lm}
\newtcbtheorem[number within=section, use counter from=thm]{prop}{Proposition}%
{
colback=white!5,
colframe=white!35!black,
fonttitle=\bfseries,
label type=proposition
}{pr}
\newtcbtheorem[number within=section, use counter from=thm]{cor}{Corollary}%
{colback=blue!5,
colframe=blue!35!black,
fonttitle=\bfseries,
label type=corollary
}{co}
\newtcbtheorem[number within=section, use counter from=thm]{axiom}{Axiom}%
{colback=black!5,
colframe=yellow!35!black,
fonttitle=\bfseries,
label type=axiom
}{ax}
\newtcbtheorem[number within=section, use counter from=thm]{defn}{Definition}%
{colback=red!5,
colframe=red!35!black,
fonttitle=\bfseries,
label type=definition
}{df}


% for <s-cr> to create \item[$\LRw$] instead of \item
\newenvironment{equivEnumerate}
 {\begin{enumerate}
	}{
\end{enumerate}}

\newtcbtheorem[number within=section, use counter from=thm]{titledBox}{}{%
  enhanced,
  sharp corners,
  colback=white,
  colframe=black,
  borderline={0.2pt}{0pt}{black},
  left=2mm,
  right=2mm,
  top=2mm,
  bottom=2mm,
  title style={opacity=1},
  attach title to upper,
  before upper={\textbf{\tcbtitletext}\quad},
  coltitle=black,
  fonttitle=\bfseries,
  separator sign={\quad}
}{box}


%better example and proof environments
\def\exampletext{Example} % If English
\colorlet{colexam}{red!55!black} % Global example color


\newtcbtheorem[number within=section, use counter from=thm]{example}{Example}{% \newtcolorbox[use counter=testexample]{testexamplebox}{%
% Example Frame Start
empty,% Empty previously set parameters
title={\exampletext \thetcbcounter #1},% use \thetcbcounter to access the testexample counter text
% Attaching a box requires an overlay
attach boxed title to top left,
   % Ensures proper line breaking in longer titles
   minipage boxed title,
% (boxed title style requires an overlay)
boxed title style={empty,size=minimal,toprule=0pt,top=4pt,left=3mm,overlay={}},
coltitle=colexam,fonttitle=\bfseries,
before=\par\medskip\noindent,parbox=false,boxsep=0pt,left=3mm,right=0mm,top=2pt,breakable,pad at break=0mm,
   before upper=\csname @totalleftmargin\endcsname0pt, % Use instead of parbox=true. This ensures parskip is inherited by box.
% Handles box when it exists on one page only
overlay unbroken={\draw[colexam,line width=.5pt] ([xshift=-0pt]title.north west) -- ([xshift=-0pt]frame.south west); },
% Handles multipage box: first page
overlay first={\draw[colexam,line width=.5pt] ([xshift=-0pt]title.north west) -- ([xshift=-0pt]frame.south west); },
% Handles multipage box: middle page
overlay middle={\draw[colexam,line width=.5pt] ([xshift=-0pt]frame.north west) -- ([xshift=-0pt]frame.south west); },
% Handles multipage box: last page
overlay last={\draw[colexam,line width=.5pt] ([xshift=-0pt]frame.north west) -- ([xshift=-0pt]frame.south west); }%
}{ex}



\newtcolorbox{ProofTcb}[1][]{%
    empty,
    title={\emph{Proof}: #1},
    attach boxed title to top left,
    minipage boxed title,
    boxed title style={empty,size=minimal,toprule=0pt,top=4pt,left=3mm,overlay={}},
    coltitle=black!55!black,
    fonttitle=\bfseries,
    before=\par\medskip\noindent,
    parbox=false,
    boxsep=0pt,
    left=3mm,
    right=0mm,
    top=2pt,
    breakable,
    pad at break=0mm,
    before upper=\csname @totalleftmargin\endcsname0pt,
    width=\dimexpr\textwidth-3mm\relax,
    overlay unbroken={\draw[black!55!black,line width=.5pt] ([xshift=-0pt]title.north west) -- ([xshift=-0pt]frame.south west); },
    overlay first={\draw[black!55!black,line width=.5pt] ([xshift=-0pt]title.north west) -- ([xshift=-0pt]frame.south west); },
    overlay middle={\draw[black!55!black,line width=.5pt] ([xshift=-0pt]frame.north west) -- ([xshift=-0pt]frame.south west); },
    overlay last={\draw[black!55!black,line width=.5pt] ([xshift=-0pt]frame.north west) -- ([xshift=-0pt]frame.south west); },
    #1
}

\newenvironment{Proof}[1][]
  {\begin{ProofTcb}[#1]}
  {\end{ProofTcb}}


%shortcuts for this file (newcommand and renewcommand)

% I like original mathcal better, but need lowercase.
\DeclareFontFamily{OT1}{pzc}{}
\DeclareFontShape{OT1}{pzc}{m}{it}{<-> s * [1.10] pzcmi7t}{}
\DeclareMathAlphabet{\mathpzc}{OT1}{pzc}{m}{it}

% Custom Commands
\newcommand{\A}{{\mathbb{A}}}
\newcommand{\F}{{\mathbb{F}}}
\newcommand{\N}{{\mathbb{N}}}
\newcommand{\Q}{{\mathbb{Q}}}
\newcommand{\R}{{\mathbb{R}}}
\newcommand{\V}{{\mathbb{V}}}
\newcommand{\Z}{{\mathbb{Z}}}
\newcommand{\C}{{\mathbb{C}}}
\renewcommand{\P}{{\mathbb{P}}}
\renewcommand{\H}{{\mathbb{H}}}

\newcommand{\CA}{{\mathcal{A}}}
\newcommand{\CB}{{\mathcal{B}}}
\newcommand{\CC}{{\mathcal{C}}}
\newcommand{\CD}{{\mathcal{D}}}
\newcommand{\CF}{{\mathcal{F}}}
\newcommand{\CG}{{\mathcal{G}}}
\newcommand{\CH}{{\mathcal{H}}}
\newcommand{\CI}{{\mathcal{I}}}
\newcommand{\CL}{{\mathcal{L}}}
\newcommand{\CN}{{\mathcal{N}}}
\newcommand{\CO}{{\mathcal{O}}}
\newcommand{\co}{{\mathpzc{o}}}
\newcommand{\CP}{{\mathbb{C}\mathrm{P}}}
\newcommand{\RP}{{\mathbb{R}\mathrm{P}}}
\newcommand{\CQ}{{\mathcal{Q}}}
\newcommand{\CR}{{\mathcal{R}}}
\newcommand{\CK}{{\mathcal{K}}}
\newcommand{\CS}{{\mathcal{S}}}
\newcommand{\CT}{{\mathcal{T}}}
\newcommand{\CV}{{\mathcal{V}}}
\newcommand{\CZ}{{\mathcal{Z}}}
\newcommand{\ff}{{\mathscr{F}}}
\newcommand{\EF}{{\EuScript{F}}}
\newcommand{\EL}{{\EuScript{L}}}
\newcommand{\EO}{{\EuScript{O}}}
\newcommand{\EC}{{\EuScript{C}}}
\newcommand{\fo}{{\mathfrak{o}}}
\newcommand{\fa}{{\mathfrak{a}}}
\newcommand{\fb}{{\mathfrak{b}}}
\newcommand{\fc}{{\mathfrak{c}}}
\newcommand{\fd}{{\mathfrak{d}}}
\newcommand{\fm}{{\mathfrak{m}}}
\newcommand{\fn}{{\mathfrak{n}}}
\newcommand{\fp}{{\mathfrak{p}}}
\newcommand{\fq}{{\mathfrak{q}}}
\newcommand{\FA}{{\mathfrak{A}}}
\newcommand{\FB}{{\mathfrak{B}}}
\newcommand{\FP}{{\mathfrak{P}}}
\newcommand{\FC}{{\mathfrak{C}}}
\newcommand{\FD}{{\mathfrak{D}}}
\newcommand{\FM}{{\mathfrak{M}}}
\newcommand{\FN}{{\mathfrak{N}}}
\newcommand{\FO}{{\mathfrak{O}}}



\newcommand{\sse}{\subseteq}
\newcommand{\nsse}{\not\subseteq}
\newcommand{\ssne}{\subsetneq}
\newcommand{\qn}{\quad\newline}
\newcommand{\placeholder}{\_\_\_\_\_}
\newcommand{\subg}{\leqslant}
%\newcommand{\subgne}{\lneqslant}
\newcommand{\supg}{\geqslant}
\newcommand{\nsubg}{\trianglelefteq}
\newcommand{\csubg}{\blacktriangleleft}
\newcommand{\nsupg}{\trianglerighteq}
\newcommand{\isosubg}{\lesssim}
\renewcommand{\mid}{\big|}
\renewcommand{\phi}{\varphi}

\newcommand{\oln}[1]{{\overline{#1}}}
\renewcommand{\bf}[1]{\textbf{#1}}

\renewcommand{\mod}{\,\operatorname{mod}\,}

\newcommand{\ev}{{\operatorname{ev}}}
\newcommand{\tr}{{\operatorname{tr}}}
\newcommand{\nm}{{\operatorname{Nm}}}
\newcommand{\op}{{\operatorname{op}}}
\newcommand{\ad}{{\operatorname{ad}}}
\newcommand{\rk}{{\operatorname{rk}}}
\newcommand{\im}{{\operatorname{im}}}
\newcommand{\id}{{\operatorname{id}}}
\newcommand{\ch}{{\operatorname{ch}}}
\newcommand{\ab}{{\operatorname{ab}}}
\newcommand{\Ch}{{\operatorname{Ch}}}
\newcommand{\Br}{{\operatorname{Br}}}
\newcommand{\Ab}{{\operatorname{Ab}}}
\newcommand{\SL}{{\operatorname{SL}}}
\newcommand{\GL}{{\operatorname{GL}}}
\newcommand{\rad}{\operatorname{rad}}
\newcommand{\sep}{\operatorname{sep}}
\newcommand{\vol}{\operatorname{vol}}
\newcommand{\ord}{\operatorname{ord}}
\newcommand{\Gal}{\operatorname{Gal}}
\newcommand{\Rep}{\operatorname{Rep}}
\newcommand{\Aut}{\operatorname{Aut}}
\newcommand{\Sym}{\operatorname{Sym}}
\newcommand{\Mod}[1][]{{}_{#1}\operatorname{Mod}}
\newcommand{\RMod}{{}_R\operatorname{Mod}}
\newcommand{\Grp}{{\textbf{Grp}}}
\newcommand{\Alg}{{\textbf{Alg}}}
\newcommand{\Set}{{\textbf{Set}}}
\newcommand{\Fld}{{\textbf{Fld}}}
\newcommand{\AlgInt}{{\textbf{AlgInt}}}
\newcommand{\fin}{{\operatorname{fin}}}
\newcommand{\Pic}{{\operatorname{Pic}}}
\newcommand{\Div}{{\operatorname{Div}}}
\newcommand{\Hom}{{\operatorname{Hom}}}
\newcommand{\Emb}{{\operatorname{Emb}}}
\newcommand{\Ext}{{\operatorname{Ext}}}
\newcommand{\End}{{\operatorname{End}}}
\newcommand{\Inn}{{\operatorname{Inn}}}
\newcommand{\Int}{{\operatorname{int}}}
\newcommand{\Irr}{{\operatorname{Irr}}}
\newcommand{\Ind}{{\operatorname{Ind}}}
\newcommand{\Out}{{\operatorname{Out}}}
\newcommand{\Jac}{{\operatorname{Jac}}}
\newcommand{\Nil}{{\operatorname{Nil}}}
\newcommand{\Tor}{{\operatorname{Tor}}}
\newcommand{\Top}{{\operatorname{Top}}}
\newcommand{\Ann}{{\operatorname{Ann}}}
\newcommand{\Syl}{{\operatorname{Syl}}}
\newcommand{\Res}{{\operatorname{Res}}}
\newcommand{\lcm}{{\operatorname{lcm}}}
\newcommand{\conj}{{\operatorname{conj}}}
\newcommand{\disc}{{\operatorname{disc}}}
\newcommand{\stab}{{\operatorname{stab}}}
\newcommand{\kdim}{{\operatorname{kdim}}}
\newcommand{\diag}{{\operatorname{diag}}}
\newcommand{\coker}{{\operatorname{coker}}}
\newcommand{\Obj}{{\operatorname{Obj}}}
\newcommand{\Mor}[1]{\operatorname{Mor}(\textbf{#1})}
\newcommand{\Frac}{{\operatorname{Frac}}}
\newcommand{\Frob}{{\operatorname{Frob}}}
\newcommand{\Spec}{{\operatorname{Spec}}}
\newcommand{\mSpec}{{\operatorname{mSpec}}}
\newcommand{\Tens}[1]{#1\otimes \_\_}
\newcommand{\Char}{{\operatorname{char}}}
\newcommand{\Span}{{\operatorname{span}}}
\newcommand{\Lim}{{\varprojlim}}
\newcommand{\colim}{\varinjlim}
\newcommand{\coLim}{{\varinjlim}}
\newcommand{\Dlim}{\lim\limits_{\longrightarrow}}
\newcommand{\coDlim}{\lim\limits_{\longleftarrow}}
\newcommand{\codim}{{\operatorname{codim}}}
\newcommand{\fHom}[1]{\operatorname{Hom}(#1, \_\_)}
\newcommand{\cofHom}[1]{\operatorname{Hom}(\_\_, #1)}
\newcommand{\trdeg}{{\operatorname{trdeg}}}
\newcommand{\acton}{\curvearrowright}
\newcommand{\actson}{\curvearrowright}

\renewcommand{\Re}{{\operatorname{Re}}}
\renewcommand{\Im}{{\operatorname{Im}}}

%for saturated ideals in the localization section
\newcommand{\LIm}{{}^{\iota} }
\newcommand{\LPre}{{}^{\pi} }

% Arrows
\newcommand{\rw}{\rightarrow}
\newcommand{\rrw}{\rightrightarrows}
\newcommand{\Rw}{\Rightarrow}
\newcommand{\lw}{\leftarrow}
\newcommand{\Lw}{\Leftarrow}
\newcommand{\lrw}{\leftrightarrow}
\newcommand{\LRw}{\Leftrightarrow}
\newcommand{\supp}{\operatorname{supp}}
\newcommand{\loc}{\operatorname{loc}}

\newcommand{\dcup}{\sqcup}
\newcommand{\Dcup}{\coprod}
\newcommand{\set}[2]{\left\{#1 \ \middle|\ #2\right\}}

\newcommand\mapsfrom{\mathrel{\reflectbox{\ensuremath{\mapsto}}}}

\newcommand{\tikzcircle}[2][red,fill=black]{\tikz[baseline=-0.5ex]\draw[#1,radius=#2] (0,0) circle ;}%

\def\rddots{\mathstrut^{.^{.^{.}}}}

\makeatletter
\newcommand*\bigcdot{\mathpalette\bigcdot@{.5}}
\newcommand*\bigcdot@[2]{\mathbin{\vcenter{\hbox{\scalebox{#2}{$\m@th#1\bullet$}}}}}
\makeatother

%% new delimiter
\DeclarePairedDelimiter{\ceil}{\lceil}{\rceil}


\setlength{\parindent}{0pt} % don't like indentation infront of paragraphs
\setlength{\parskip}{1ex}   %so that there is still space between paragarphs


%for Dynkin diagrams
\def\row#1/#2!{#1_{\IfStrEq{#2}{}{n}{#2}} & \dynkin{#1}{#2}\\}
\newcommand{\tble}[1]{
   \renewcommand*\do[1]{\row##1!}
   \[
      \begin{array}{ll}\docsvlist{#1}\end{array}
   \]
}

\makeatletter
\renewcommand*\env@matrix[1][*\c@MaxMatrixCols c]{%
  \hskip -\arraycolsep
  \let\@ifnextchar\new@ifnextchar
  \array{#1}}
\makeatother

\usepackage{import}
\usepackage{pdfpages}
\usepackage{transparent}

\newcommand{\incfig}[2][1]{%
    \def\svgwidth{#1\columnwidth}
    \import{./figures/}{#2.pdf_tex}
}

\pdfsuppresswarningpagegroup=1


\definecolor{LinkColour}{HTML}{A64A3B} % favourite
% \definecolor{LinkColour}{HTML}{8C3D32} % 2nd place
\hypersetup{
  colorlinks,
  citecolor=black,
  filecolor=magenta,
  linkcolor=LinkColour,
  urlcolor=blue,
}

\makeindex

\begin{document}
\titleUL
\tableofcontents

\subsection*{Summary}

It as assumed that the reader is comfortable with the concept introduced in a first course in analysis and
linear algebra, importantly that smooth functions and vector-spaces are comfortable concepts.

\chapter*{Introduction Differential Equations}

To avoid super-script counting, the book shall work with smooth functions (or else the order of
differentiability must be kept track). This shall not change the theory for ODE's and PDEs, especially since
later it shall be shown that even working with non-smooth differentiable functions shall end up having the
same theory (due to the advent of Distribution Theory \footnote{see chapter ref:HERE})

Let $f : \R \to \R$. When working smooth functions, we may take the derivative
of the function. This gives the following function taking smooth functions and returning smooth functions:

\[
	f \mapsto f'
\]
Due to the fundamental theorem of calculus, the inverse of this procedure is called \emph{integration}, and gives a unique inverse up to a constant:
\[
	f \mapsto F + c \qquad c \in \R
\]
where $F$ may be defined as $F(x) = \int_t^x f(x)dx$ for appropriate $t$. Using this result, it is possible
to find the smooth function that satisfies an arbitrary degree of differentiation, for example:
\[
	f^{(n)} = 0
\]
which is a polynomial of degree $n$, or more generally
\[
	f^{(n)} = g \qquad \text{equiv.} \qquad f^{(n)} - g = 0
\]
for some smooth functions $g$\footnote{Note that the solution may not always come in an elementary form, that
is expressible as a polynomial-combination and composition of $x$, $\log$, $e^x$, or a trigonometric
identity}.


Given this simple scenario, the question may arise if we may solve more
complicated algebraic expression. For example, does there exist a solution to any of these:
\[
	f = f'\qquad f' + f = 0\qquad f'' + f' + f = 0
\]
More precisely, we ask if any arbitrary algebraic expression on a smooth function $f$ combined by using the 4 usual arithmetic operators
($+,-,\times,\div$) along with differentiation yields a solution:
\[
	p(f) = 0
\]
more commonly written as:
\[
	F(x, f(x), f'(x), ..., f^{(n)}(x)) = 0
\]

The study of the solution to such equations is called the study of \emph{Ordinary Differential Equations}
(ODEs). If we have a smooth function of the form $f: \R \to \R^m$, this is still the study of ODEs,
where it is solving a \emph{system of ODEs}. When working with multi-variable functions and partial derivatives, the study of the solution to such
equations is called the study of \emph{Partial Differential Equations} (PDEs). The study of ODEs can be
thought of as a special case of the  study of PDEs.

This book shall start with the study of ODEs in the first chapter, and proceed with the rest of the chapter
with the study of PDEs. Only one chapter is dedicated to ODEs as they are a solved problem where most ODEs
have a general solution that can be found algorithmically, while this is provably not the case for most PDEs.



From a more physical perspective, differential equations will give you how certain parameters relate to each,
for PDEs it is usually over time. As a simple example, take:
\begin{equation}\label{eq:pdeIntuition}
  \frac{\partial }{\partial t} = \frac{\partial }{\partial x} \qquad \text{or} \qquad \frac{\partial }{\partial t}
  - \frac{\partial }{\partial x} = 0
\end{equation}
What this equation is giving us is a restriction we require a function $f(t,x)$ to obey, that is:
\[
	\frac{\partial f}{\partial t} - \frac{\partial f}{\partial x} = 0
\]
For example $f(x,t) = x+t$ would satisfy this equation. One key skill to learn when solving problems in PDEs is to be
able to ``read'' such an equation. For example, equation~(\ref{eq:pdeIntuition}) tells us that the speed at which our
``system'' is evolving over time is identically propositional to the speed at which our ``system'' is moving. Sometimes,
it is quite subtle what an equation may represent. For example:
\[
	\frac{\partial^2}{\partial t^2} - \frac{\partial^2}{\partial x^2} = 0
\]
is very similar to equation~(\ref{eq:pdeIntuition}), however it represents how a \emph{wave} evolves over
time, with behavior that is very different from equation~(\ref{eq:pdeIntuition}).


The reason why finding solutions to PDEs is interesting to me is
because the derivative can give lot's of local information about a system, and sometimes that's all the information we
have. For example, it is much easier to describe the force of gravity as

\[
	F = ma
\]
where $a = dv/dt$, and $v = d(\textrm{displacement})/dt = dD/dt$. We can use this equation to construct a gravitational
field to how any object at any point will have a certain force on them represented by a vector pointing downwards times
the mass of the object, but this ``global solution'' has a lot more information than is needed. More generally, it is
often easier to come up with the ``localizations'', since they are easier to define or control most of the time.

Also, after seeing how heat equations and diffusion equations are described I have a new motivation:
I wonder if at some point these techniques of relating ``local information'' in describing how it's dynamics evolve (in
physics, it would be how it changes in time, in more general scenarios it can be a recursive function or a sequence)
will at some point lend itself useful in solving more general problems. For example, let's say we have a particularly complicated
mathematical object. Let's say we can solve it in some nice special cases, perhaps with some restricting assumptions.
Can we then step-by-step evolve the problem (or model an evolution of the problem) to eliminate the restrictions? Can we
localize in some way to simplify it, and glue that information together? These feel like modelling something like a PDE.

As a final motivator, every differential form has a ``canonical representative'' in that there always exists
representative $\omega$ such that $\omega$ vanishes under the Laplacian: $\Delta \omega = 0$ (functions, and more
generally forms, that vanishes under the laplacian are called \emph{harmonic}). This is the beginning of \emph{Hodge
Theory}, which is central in studying the cohomology groups of smooth manifolds. In the same vein as geometry but from
another perspective, PDEs are also used to model  \emph{Ricci flow}, which is about the flow of a manifold given certain
parameters as it changes across time.

\subsection*{Ordinary Differential Equations}



We focus on solving \emph{Ordinary differential equations} with the hope that we may generalize to the non-linear case. For example:

\[
	5y'' + 1y' -3y =0
\]

is a linear differential equation. Since it is equal to $0$, it is called a \emph{homogeneous} linear ODE. Since
trigonometric functions have very simple derivative properties, they are also studied regularly. Though this might seem
limiting, these in fact cover a satisfyingly large sample of scenario's which we may encounter in physics. If your
motivation is within pure mathematics, then having a standard repertoire to solve linear and trigonometric ODE's allows
us to use these as ``local tools'' to solve non-linear ODE's or ODE's approximated by a trigonometric
sequence\footnote{If you know some more advanced analysis, think of functions being equal to a series of
trigonometric functions or polynomials}. Here is small a list of types of equations and techniques used to solve ODE's:

\begin{enumerate}
	\item First order ODE:
	\begin{enumerate}
		 \item Separable: if the equation can be written in the form: $\frac{dy(x)}{dx} = \frac{m(x)}{n(y)}$, or equivalently:
			  	\[
			  		\frac{dy(x)}{dx} + \frac{m(x)}{n(y)} = 0
			  	\]
			  	Then we can solve this ODE
			  \item Linear (Using an integral factor): If the equation can be written in the form:
			  	\[
			  		\frac{dy(x)}{dx} + p(x)y(x) = g(x)
			  	\]
			  	then we can solve this ODE
	\end{enumerate}
	\item Second order ODE:
		\begin{enumerate}
			\item If the ODE has constant coefficients, is linear, and homogeneous:
				\[
					ay'' + by' + cy  =0
			\]
			then we can solve this ode, with $y(x) = Ae^{rx}$ for some $r \in \C$ and $A \in \C$.
		\item If the ODE is Equidimensional (Cauchy-Euler), linear, homogeneous:
			\[
				ax^2y'' + bxy' + cy  =0
			\]
			then the solution is of the form $y(x) = Ax^p$ for some constant $p \in \C$.
		\end{enumerate}
\end{enumerate}

There are many more that I've not listed here, I just wanted to give a list to give an idea of what was done.

\chapter{First-Order Scalar Equations}\label{cha:firstOrderODE}

\subsection*{Summary}
In this chapter, we establish the fundamental theory for scalar first-order ordinary differential equations. We begin by defining the geometric interpretation of an ODE as a vector field, establishing the conditions required for the existence and uniqueness of solutions (Picard-Lindelöf), and then proceed to classify and solve integrable forms: separable, linear, and exact equations.

\section{The General Initial Value Problem (IVP)}

We begin our study with the scalar first-order ordinary differential equation in the normal form:
\begin{equation}\label{eq:normalFormODE}
    y' = f(t, y)
\end{equation}
where $y(t)$ is the unknown function and $f: U \to \R$ is a function defined on some open domain $U \sse \R^2$. Physically, if we view $t$ as time and $y$ as position, $f$ prescribes the velocity of a particle based on its current time and location. Geometrically, $f$ assigns a slope to every point in $U$, generating a \emph{slope field} (or direction field). A solution $\phi(t)$ is a curve whose tangent at every point $(t, \phi(t))$ matches the slope prescribed by $f$.

\begin{defn}{Initial Value Problem (IVP)}{ivpDef}
    An \emph{Initial Value Problem} consists of an ODE combined with a specified value at a point $t_0$:
    \[
        \begin{cases}
            y' = f(t, y) \\
            y(t_0) = y_0
        \end{cases}
    \]
    A solution to the IVP on an interval $I$ containing $t_0$ is a differentiable function $\phi: I \to \R$ such that $\phi'(t) = f(t, \phi(t))$ for all $t \in I$ and $\phi(t_0) = y_0$.
\end{defn}

Before we attempt to solve these equations, we must address the fundamental questions of analysis: Does a solution exist? Is it unique?

\begin{thm}{Picard-Lindelöf Existence and Uniqueness}{picardLindelof}
    Let $D = \set{(t,y) \in \R^2}{|t - t_0| \le a, |y - y_0| \le b}$ be a rectangular region centered at the initial data. Suppose that:
    \begin{enumerate}
        \item $f(t,y)$ is continuous on $D$.
        \item $f$ is \emph{Lipschitz continuous} with respect to $y$ on $D$. That is, there exists a constant $L > 0$ such that:
        \[
            |f(t, y_1) - f(t, y_2)| \le L|y_1 - y_2| \qquad \forall (t, y_1), (t, y_2) \in D
        \]
    \end{enumerate}
    Then there exists a unique solution $y(t)$ to the IVP defined on some interval $|t - t_0| \le h$ (where $h \le a$).
\end{thm}

\begin{Proof}
    The proof relies on converting the ODE into an integral equation:
    \[
        y(t) = y_0 + \int_{t_0}^t f(s, y(s)) \, ds
    \]
    and defining the Picard Iteration operator $T: y \mapsto y_0 + \int f(s, y(s))$. Under the Lipschitz condition, one proves that $T$ is a contraction mapping on a suitable Banach space of continuous functions. By the Banach Fixed Point Theorem, a unique fixed point (solution) exists.
\end{Proof}

Note that if $\frac{\partial f}{\partial y}$ exists and is continuous, the Lipschitz condition is automatically satisfied locally.

\section{Separable Equations}\label{sec:separable}

The simplest class of first-order equations are those where the dependence on $t$ and $y$ can be decoupled.

\begin{defn}{Separable Equation}{separableDef}
    A first-order ODE is \emph{separable} if it can be written in the form:
    \[
        \frac{dy}{dt} = g(t)h(y)
    \]
    or equivalently, if $h(y) \ne 0$:
    \begin{equation}\label{eq:separatedForm}
        \frac{1}{h(y)} \, dy = g(t) \, dt
    \end{equation}
\end{defn}

The solution technique relies on the chain rule. Integrating \eqref{eq:separatedForm} with respect to $t$:
\[
    \int \frac{1}{h(y(t))} \frac{dy}{dt} \, dt = \int g(t) \, dt \implies \int \frac{1}{h(y)} \, dy = \int g(t) \, dt
\]
This yields an implicit relationship $H(y) = G(t) + C$.

\begin{example}{Solving a Separable IVP}{sepEx1}
    Solve the IVP:
    \[
        y' = -2ty^2, \qquad y(0) = 1
    \]
    \emph{Solution:}
    We separate variables (assuming $y \ne 0$):
    \[
        \frac{dy}{y^2} = -2t \, dt
    \]
    Integrating both sides:
    \[
        \int y^{-2} \, dy = \int -2t \, dt \implies -\frac{1}{y} = -t^2 + C \implies y(t) = \frac{1}{t^2 - C}
    \]
    Using the initial condition $y(0) = 1$:
    \[
        1 = \frac{1}{0 - C} \implies C = -1
    \]
    Thus, the solution is $y(t) = \frac{1}{t^2 + 1}$. Note that this solution is defined for all $t \in \R$. Had the initial condition been $y(0) = -1$, we would have found $y(t) = \frac{1}{t^2 - 1}$, which blows up at $t = \pm 1$, illustrating that solutions to non-linear ODEs may exist only locally.
\end{example}

\textbf{Remark on Singular Solutions:} In the separation step, we divided by $h(y)$. The roots of $h(y) = 0$ represent constant solutions (equilibrium solutions) $y(t) = c$. These must be checked separately to see if they satisfy the initial conditions or if solutions converge to them.

\section{Linear Equations and Integrating Factors}

Linear equations are the only class of ODEs for which a complete, closed-form general solution is always guaranteed (provided the coefficient functions are continuous).

\begin{defn}{First-Order Linear ODE}{linearDef}
    A first-order ODE is \emph{linear} if it can be written as:
    \[
        y' + p(t)y = g(t)
    \]
    If $g(t) \equiv 0$, the equation is \emph{homogeneous}; otherwise, it is \emph{inhomogeneous}.
\end{defn}

To solve this, we seek a function $\mu(t)$, called an \emph{integrating factor}, such that multiplying the equation by $\mu(t)$ transforms the left-hand side into a total derivative of a product.
\[
    \mu(t)y' + \mu(t)p(t)y = \mu(t)g(t)
\]
We want the LHS to be $\frac{d}{dt}[\mu(t)y] = \mu(t)y' + \mu'(t)y$. Matching terms, we require:
\[
    \mu'(t) = \mu(t)p(t)
\]
This is a separable equation for $\mu$: $\frac{d\mu}{\mu} = p(t)dt$, yielding $\mu(t) = \exp\left( \int p(t) \, dt \right)$.

\begin{thm}{General Solution of Linear Equations}{linearSolThm}
    The general solution to $y' + p(t)y = g(t)$ is given by:
    \[
        y(t) = \frac{1}{\mu(t)} \left[ \int \mu(t)g(t) \, dt + C \right]
    \]
    where $\mu(t) = e^{\int p(t) \, dt}$.
\end{thm}

\begin{Proof}
    Multiply by $\mu(t)$. Then $\frac{d}{dt}[\mu(t)y] = \mu(t)g(t)$. Integrating both sides yields $\mu(t)y = \int \mu(t)g(t) dt + C$. Dividing by $\mu(t)$ gives the result.
\end{Proof}

\begin{example}{Linear Equation with Discontinuous Forcing}{linearEx}
    Consider $y' + 2y = g(t)$ where $y(0)=0$ and $g(t) = \chi_{[0,1]}$ (the characteristic function of $[0,1]$).

    \emph{Solution:}
    The integrating factor is $\mu(t) = e^{\int 2 dt} = e^{2t}$.
    \[
        \frac{d}{dt}[e^{2t}y] = e^{2t}g(t)
    \]
    Integrating from $0$ to $t$:
    \[
        e^{2t}y(t) - e^0y(0) = \int_0^t e^{2s}g(s) \, ds
    \]
    If $0 \le t \le 1$:
    \[
        e^{2t}y(t) = \int_0^t e^{2s}(1) \, ds = \frac{1}{2}(e^{2t}-1) \implies y(t) = \frac{1}{2}(1 - e^{-2t})
    \]
    If $t > 1$:
    \[
        e^{2t}y(t) = \int_0^1 e^{2s} \, ds + \int_1^t 0 \, ds = \frac{1}{2}(e^2 - 1) \implies y(t) = \frac{1}{2}(e^2 - 1)e^{-2t}
    \]
    Notice that while $g(t)$ is discontinuous, the solution $y(t)$ is continuous, though $y'(t)$ is discontinuous at $t=1$. This ``smoothing'' property is characteristic of ODEs.
\end{example}

\section{Exact Equations and Differential Forms}

Here we take a slightly more sophisticated view. We can rewrite a first-order equation $y' = f(x,y)$ symmetrically using differential forms:
\[
    M(x,y) \, dx + N(x,y) \, dy = 0
\]
We ask: does there exist a scalar function (a potential) $\psi(x,y)$ such that $d\psi = M dx + N dy$? If so, the differential equation becomes $d\psi = 0$, implying the general solution is the level set $\psi(x,y) = C$.

\begin{defn}{Exact Differential Equation}{exactDef}
    The equation $M(x,y)dx + N(x,y)dy = 0$ is \emph{exact} in a region $R$ if there exists a function $\psi$ such that:
    \[
        \frac{\partial \psi}{\partial x} = M \qquad \text{and} \qquad \frac{\partial \psi}{\partial y} = N
    \]
\end{defn}

By equality of mixed partials ($\psi_{xy} = \psi_{yx}$), a necessary condition for exactness is:
\begin{equation}\label{eq:exactCondition}
    \frac{\partial M}{\partial y} = \frac{\partial N}{\partial x}
\end{equation}
On a simply connected domain (like a rectangle or all of $\R^2$), the Poincaré Lemma states this condition is also sufficient.

\begin{example}{Solving an Exact Equation}{exactEx}
    Solve $(2xy + 3)dx + (x^2 - 1)dy = 0$.

    \emph{Check Exactness:}
    \[
        \partial_y (2xy + 3) = 2x \qquad \partial_x (x^2 - 1) = 2x
    \]
    The equation is exact. We seek $\psi$ such that $\psi_x = 2xy + 3$. Integrating with respect to $x$:
    \[
        \psi(x,y) = x^2y + 3x + h(y)
    \]
    Now differentiate with respect to $y$ and match $N$:
    \[
        \psi_y = x^2 + h'(y) \overset{!}{=} x^2 - 1 \implies h'(y) = -1 \implies h(y) = -y
    \]
    Thus $\psi(x,y) = x^2y + 3x - y$. The implicit solution is $x^2y + 3x - y = C$.
\end{example}

\subsection*{Integrating Factors for Non-Exact Equations}
If $M dx + N dy = 0$ is not exact, we may look for a function $\mu(x,y)$ such that $(\mu M) dx + (\mu N) dy = 0$ is exact. This requires:
\[
    (\mu M)_y = (\mu N)_x \implies \mu_y M + \mu M_y = \mu_x N + \mu N_x
\]
This is a PDE for $\mu$, which is generally harder to solve than the original ODE! However, if we assume $\mu$ depends only on $x$ (i.e., $\mu_y = 0$), the condition simplifies to:
\[
    \frac{d\mu}{dx} N = \mu (M_y - N_x) \implies \frac{1}{\mu}\frac{d\mu}{dx} = \frac{M_y - N_x}{N}
\]
If the RHS depends only on $x$, we can solve for $\mu$. A similar check can be done assuming $\mu = \mu(y)$.

\section{Bernoulli Equations}

We conclude this chapter with a classic non-linear equation that can be linearized. A \emph{Bernoulli equation} has the form:
\begin{equation}\label{eq:bernoulli}
    y' + p(t)y = g(t)y^n \qquad n \in \R, n \ne 0, 1
\end{equation}
The substitution $v = y^{1-n}$ linearizes this equation.
\[
    v' = (1-n)y^{-n}y'
\]
Multiplying \eqref{eq:bernoulli} by $(1-n)y^{-n}$:
\[
    (1-n)y^{-n}y' + (1-n)p(t)y^{1-n} = (1-n)g(t)
\]
\[
    v' + (1-n)p(t)v = (1-n)g(t)
\]
This is now a standard linear equation for $v(t)$, which can be solved using an integrating factor.


\chapter{Second-Order Linear Equations}\label{cha:secondOrderODE}

\subsection*{Summary}
We now move to second-order equations. We focus on \emph{linear} equations because their solution space forms a vector space, allowing us to use tools from linear algebra. We will rigorously define linear independence using the Wronskian, prove Abel's Identity, solve constant coefficient equations via characteristic polynomials, and derive the Variation of Parameters formula for inhomogeneous equations.

\section{Definitions and Theoretical Framework}

We begin by establishing the precise vocabulary used to classify these equations.

\begin{defn}{Second-Order Linear ODE}{secondOrderLinearDef}
    A second-order ordinary differential equation is called \emph{linear} if it can be written in the form:
    \begin{equation}\label{eq:generalLinear2nd}
        P(t)y'' + Q(t)y' + R(t)y = G(t)
    \end{equation}
    where $P, Q, R, G$ are continuous functions on an interval $I$, and $P(t) \ne 0$ on $I$.

    If we divide by $P(t)$, we obtain the \emph{standard form}:
    \[
        y'' + p(t)y' + q(t)y = g(t)
    \]
\end{defn}

\begin{defn}{Homogeneous vs. Inhomogeneous}{homogInhomogDef}
    Referring to the standard form $y'' + p(t)y' + q(t)y = g(t)$:
    \begin{itemize}
        \item If $g(t) \equiv 0$ for all $t$, the equation is called \emph{homogeneous}.
        \item If $g(t) \not\equiv 0$, the equation is called \emph{inhomogeneous} (or non-homogeneous).
    \end{itemize}
\end{defn}

The distinction is crucial: homogeneous equations describe the ``natural'' behavior of a system, while
inhomogeneous equations describe the system's response to an external ``forcing'' function $g(t)$. The
following theorem guarantees that as long as the coefficient functions $p(t)$ and $q(t)$ are continuous, the
solution extends across the entire interval of continuity.

\begin{thm}{Existence and Uniqueness for Linear Equations}{linearExistUnique}
    Consider the Initial Value Problem (IVP):
    \[
        y'' + p(t)y' + q(t)y = g(t), \qquad y(t_0) = y_0, \quad y'(t_0) = y'_0
    \]
    If $p(t)$, $q(t)$, and $g(t)$ are continuous on an open interval $I$ containing $t_0$, then there exists a unique solution $y(t)$ defined on the entire interval $I$.
\end{thm}

\begin{Proof}
    This is a special case of the general Existence and Uniqueness theorem for linear systems (which we will cover in Chapter \ref{cha:systems}). We convert the second-order scalar equation into a first-order system by setting $x_1 = y$ and $x_2 = y'$. This yields $\vec{x}' = A(t)\vec{x} + \vec{g}(t)$. Since the entries of $A(t)$ (which are $p(t)$ and $q(t)$) are continuous on $I$, the Picard-Lindelöf theorem for systems guarantees a unique global solution on $I$.
\end{Proof}

\begin{example}{Interval of Validity}{intervalValidityEx}
    Determine the largest interval on which the IVP has a unique solution:
    \[
        (t^2 - 4)y'' + 3ty' + (\cos t)y = e^t, \qquad y(0) = 1, \quad y'(0) = 0
    \]
    \emph{Solution:}
    First, write the equation in standard form by dividing by $(t^2 - 4)$:
    \[
        y'' + \frac{3t}{t^2 - 4}y' + \frac{\cos t}{t^2 - 4}y = \frac{e^t}{t^2 - 4}
    \]
    The coefficients are discontinuous at $t = \pm 2$. The initial point is $t_0 = 0$. The largest interval containing $0$ where all functions are continuous is $(-2, 2)$. Thus, the solution is guaranteed to exist uniquely on $t \in (-2, 2)$.
\end{example}

\section{The Structure of Solutions}

The set of solutions to a homogeneous linear ODE forms a vector space.

\begin{thm}{Principle of Superposition}{superposition}
    If $y_1$ and $y_2$ are solutions to the homogeneous equation $L[y] = y'' + p(t)y' + q(t)y = 0$, then any linear combination $c_1 y_1 + c_2 y_2$ is also a solution.
\end{thm}

\begin{Proof}
    This follows from the linearity of the differential operator $L$.
    \[
        L[c_1 y_1 + c_2 y_2] = c_1 L[y_1] + c_2 L[y_2] = c_1(0) + c_2(0) = 0
    \]
\end{Proof}

To describe \emph{all} solutions, we need a ``basis'' for this solution space.

\begin{defn}{Wronskian}{wronskianDef}
    Let $y_1, y_2$ be two differentiable functions. The \emph{Wronskian} of $y_1$ and $y_2$ is the determinant:
    \[
        W[y_1, y_2](t) = \det \begin{pmatrix} y_1(t) & y_2(t) \\ y'_1(t) & y'_2(t) \end{pmatrix} = y_1(t)y'_2(t) - y_2(t)y'_1(t)
    \]
\end{defn}

The Wronskian is crucial because it allows us to test if solutions are linearly independent at a specific point, which is sufficient for linear independence on the entire interval.

\begin{thm}{Abel's Identity}{abelIdentity}
    For the ODE $y'' + p(t)y' + q(t)y = 0$, the Wronskian $W(t)$ satisfies the first-order differential equation $W' = -p(t)W$. Thus,
    \[
        W(t) = C e^{-\int p(t)dt}
    \]
    where $C$ is a constant determined by the initial conditions.
\end{thm}

\begin{Proof}
    Compute the derivative of the Wronskian definition:
    \begin{align*}
        W' &= (y_1 y_2' - y_2 y_1')' \\
        &= (y_1' y_2' + y_1 y_2'') - (y_2' y_1' + y_2 y_1'') \\
        &= y_1 y_2'' - y_2 y_1''
    \end{align*}
    Since $y_1$ and $y_2$ satisfy the ODE $y'' = -p(t)y' - q(t)y$, substitute this into the expression for $W'$:
    \begin{align*}
        W' &= y_1[-p(t)y_2' - q(t)y_2] - y_2[-p(t)y_1' - q(t)y_1] \\
        &= -p(t)(y_1 y_2' - y_2 y_1') - q(t)(y_1 y_2 - y_2 y_1) \\
        &= -p(t)W - q(t)(0) \\
        &= -p(t)W
    \end{align*}
    This is a separable first-order linear equation, the solution of which is $W(t) = Ce^{-\int p(t)dt}$.
\end{Proof}

\begin{thm}{General Solution of Homogeneous Equations}{genSolHomog}
    Let $y_1, y_2$ be two solutions to the homogeneous equation $y'' + p(t)y' + q(t)y = 0$ on an interval $I$. If there exists a point $t_0 \in I$ such that $W[y_1, y_2](t_0) \ne 0$, then $\{y_1, y_2\}$ form a \emph{Fundamental Set of Solutions}.

    The general solution is:
    \[
        y(t) = c_1 y_1(t) + c_2 y_2(t)
    \]
\end{thm}

\begin{Proof}
    Let $y(t)$ be any solution. We want to find $c_1, c_2$ such that $y(t) = c_1 y_1(t) + c_2 y_2(t)$. Consider the initial conditions at $t_0$: $y(t_0) = Y_0$ and $y'(t_0) = Y_0'$. We set up the linear system:
    \[
        \begin{cases}
            c_1 y_1(t_0) + c_2 y_2(t_0) = Y_0 \\
            c_1 y_1'(t_0) + c_2 y_2'(t_0) = Y_0'
        \end{cases}
    \]
    The determinant of the coefficient matrix is exactly $W[y_1, y_2](t_0)$. Since $W(t_0) \ne 0$, a unique pair $(c_1, c_2)$ exists. By the Uniqueness Theorem, since $c_1 y_1 + c_2 y_2$ satisfies the ODE and the same initial conditions as $y(t)$, they must be the same function.
\end{Proof}

\begin{example}{Checking Linear Independence}{wronskianEx}
    Show that $y_1(t) = e^t$ and $y_2(t) = e^{-t}$ form a fundamental set of solutions for $y'' - y = 0$.

    \emph{Solution:}
    Check if they solve the ODE:
    \[
        y_1'' - y_1 = (e^t) - e^t = 0 \quad \checkmark \qquad y_2'' - y_2 = (e^{-t}) - e^{-t} = 0 \quad \checkmark
    \]
    Now compute the Wronskian:
    \[
        W(t) = \det \begin{pmatrix} e^t & e^{-t} \\ e^t & -e^{-t} \end{pmatrix} = (e^t)(-e^{-t}) - (e^{-t})(e^t) = -1 - 1 = -2
    \]
    Since $W(t) = -2 \ne 0$, the functions are linearly independent.
\end{example}

\section{Homogeneous Equations with Constant Coefficients}

We now turn to solving specific equations. The simplest case is when $P, Q, R$ are constants.
\begin{equation}\label{eq:constCoeff}
    ay'' + by' + cy = 0 \qquad a,b,c \in \R, a \ne 0
\end{equation}
Motivated by the property that exponentials reproduce themselves under differentiation, we guess a solution of the form $y = e^{rt}$. Substituting this into \eqref{eq:constCoeff}:
\[
    a(r^2 e^{rt}) + b(r e^{rt}) + c(e^{rt}) = 0 \implies e^{rt}(ar^2 + br + c) = 0
\]
Since $e^{rt} \ne 0$, we must solve the \emph{Characteristic Equation}:
\begin{equation}\label{eq:charEq}
    ar^2 + br + c = 0
\end{equation}
The roots are $r = \frac{-b \pm \sqrt{b^2 - 4ac}}{2a}$.

\begin{thm}{Solutions to Constant Coefficient Equations}{constCoeffCases}
    Depending on the discriminant $\Delta = b^2 - 4ac$, we have three cases for the general solution $y(t)$:
    \begin{enumerate}
        \item \emph{Distinct Real Roots ($\Delta > 0$):} Two roots $r_1 \ne r_2$.
        \[ y(t) = c_1 e^{r_1 t} + c_2 e^{r_2 t} \]
        \item \emph{Repeated Real Root ($\Delta = 0$):} One root $r = -b/2a$.
        \[ y(t) = c_1 e^{rt} + c_2 t e^{rt} \]
        \item \emph{Complex Conjugate Roots ($\Delta < 0$):} Roots $r = \lambda \pm i\mu$.
        \[ y(t) = e^{\lambda t} (c_1 \cos(\mu t) + c_2 \sin(\mu t)) \]
    \end{enumerate}
\end{thm}

\begin{Proof}
    Case 1 follows directly from the characteristic equation yielding distinct $r_1, r_2$, and the Wronskian $W[e^{r_1 t}, e^{r_2 t}] = (r_2 - r_1)e^{(r_1+r_2)t} \ne 0$.

    Case 2 is proven via reduction of order. Let $y_1 = e^{rt}$. Guess $y_2(t) = v(t)e^{rt}$. Then:
    \[
        y_2' = (v' + rv)e^{rt}, \quad y_2'' = (v'' + 2rv' + r^2v)e^{rt}
    \]
    Substituting into the ODE $ay'' + by' + cy = 0$:
    \[
        a(v'' + 2rv' + r^2v) + b(v' + rv) + cv = 0
    \]
    Rearranging terms by derivatives of $v$:
    \[
        av'' + (2ar + b)v' + (ar^2 + br + c)v = 0
    \]
    Since $r$ satisfies the characteristic equation, the coefficient of $v$ is zero. Since $r = -b/2a$ is a repeated root, $2ar + b = 2a(-b/2a) + b = 0$, so the coefficient of $v'$ is also zero. We are left with $av'' = 0 \implies v(t) = c_1 t + c_2$. Choosing the simplest non-constant solution gives $v(t) = t$, so $y_2 = te^{rt}$.

    Case 3 uses Euler's formula $e^{(\lambda+i\mu)t} = e^{\lambda t}(\cos \mu t + i \sin \mu t)$. Since the ODE is linear, the real and imaginary parts of a complex solution are also solutions.
\end{Proof}

\begin{example}{Solving Constant Coefficient Problems}{constCoeffEx}
    Find the general solution for:
    \begin{enumerate}
        \item $y'' - 5y' + 6y = 0$
        \item $y'' + 4y' + 4y = 0$
        \item $y'' + 4y' + 13y = 0$
    \end{enumerate}
    \emph{Solutions:}
    \begin{enumerate}
        \item Char. Eq: $r^2 - 5r + 6 = (r-2)(r-3) = 0$. Roots: $r_1=2, r_2=3$.
        \[ y(t) = c_1 e^{2t} + c_2 e^{3t} \]
        \item Char. Eq: $r^2 + 4r + 4 = (r+2)^2 = 0$. Root: $r=-2$ (repeated).
        \[ y(t) = c_1 e^{-2t} + c_2 t e^{-2t} \]
        \item Char. Eq: $r^2 + 4r + 13 = 0$. Using quadratic formula:
        \[ r = \frac{-4 \pm \sqrt{16 - 52}}{2} = -2 \pm 3i \]
        Here $\lambda = -2, \mu = 3$.
        \[ y(t) = e^{-2t}(c_1 \cos(3t) + c_2 \sin(3t)) \]
    \end{enumerate}
\end{example}

\section{Cauchy-Euler Equations}

\begin{defn}{Cauchy-Euler Equation}{cauchyEulerDef}
    An \emph{Equidimensional} or \emph{Cauchy-Euler} equation is of the form:
    \[
        at^2 y'' + bty' + cy = 0 \qquad t > 0
    \]
\end{defn}

We guess solutions of the form $y = t^m$. Substituting:
\[
    at^2(m(m-1)t^{m-2}) + bt(mt^{m-1}) + c(t^m) = 0
\]
This simplifies to the auxiliary equation:
\[
    am(m-1) + bm + c = 0 \implies am^2 + (b-a)m + c = 0
\]

\begin{thm}{Solutions to Cauchy-Euler Equations}{cauchyEulerThm}
    Let $m_1, m_2$ be the roots of $am^2 + (b-a)m + c = 0$.
    \begin{enumerate}
        \item \emph{Distinct Real Roots:} $y(t) = c_1 t^{m_1} + c_2 t^{m_2}$
        \item \emph{Repeated Real Root ($m_1=m_2=m$):} $y(t) = c_1 t^m + c_2 t^m \ln(t)$
        \item \emph{Complex Roots ($m = \lambda \pm i\mu$):}
        \[ y(t) = t^\lambda \left[ c_1 \cos(\mu \ln t) + c_2 \sin(\mu \ln t) \right] \]
    \end{enumerate}
\end{thm}

\begin{Proof}
    Substitute $t = e^x$ (so $x = \ln t$). Then use the chain rule:
    \[
        \frac{dy}{dt} = \frac{dy}{dx}\frac{dx}{dt} = \frac{1}{t}\frac{dy}{dx}
    \]
    \[
        \frac{d^2y}{dt^2} = \frac{d}{dt}\left(\frac{1}{t}\frac{dy}{dx}\right) = -\frac{1}{t^2}\frac{dy}{dx} + \frac{1}{t}\frac{d^2y}{dx^2}\frac{1}{t} = \frac{1}{t^2}\left(\frac{d^2y}{dx^2} - \frac{dy}{dx}\right)
    \]
    Substituting these into the original ODE yields a constant coefficient equation in $x$:
    \[
        a\left(\frac{d^2y}{dx^2} - \frac{dy}{dx}\right) + b\frac{dy}{dx} + cy = 0 \implies a\frac{d^2y}{dx^2} + (b-a)\frac{dy}{dx} + cy = 0
    \]
    The auxiliary equation for $t$ is exactly the characteristic equation for $x$. The solutions in terms of $x$ (e.g., $e^{m_1 x}$) become solutions in terms of $t$ (e.g., $(e^x)^{m_1} = t^{m_1}$).
\end{Proof}

\begin{example}{Cauchy-Euler Application}{cauchyEx}
    Solve $t^2 y'' - 3ty' + 3y = 0$ for $t>0$.

    \emph{Solution:}
    Auxiliary equation: $1m(m-1) - 3m + 3 = 0 \implies m^2 - 4m + 3 = 0$.
    Factors to $(m-1)(m-3)=0$. Roots are $1, 3$.
    \[
        y(t) = c_1 t + c_2 t^3
    \]
\end{example}

\section{Inhomogeneous Equations}

To solve $y'' + p(t)y' + q(t)y = g(t)$, we rely on the following structural theorem.

\begin{thm}{General Solution of Inhomogeneous Equations}{inhomogStruct}
    The general solution to the inhomogeneous equation is:
    \[
        y(t) = y_h(t) + y_p(t)
    \]
    where $y_h(t)$ is the general homogeneous solution and $y_p(t)$ is \emph{any} particular solution.
\end{thm}

\begin{Proof}
    Let $L[y] = y'' + p(t)y' + q(t)y$. Let $Y$ be any solution such that $L[Y] = g$. Let $y_p$ be a particular solution such that $L[y_p] = g$.
    Then $L[Y - y_p] = L[Y] - L[y_p] = g - g = 0$. Thus $Y - y_p$ solves the homogeneous equation, so $Y - y_p = y_h$. Hence $Y = y_h + y_p$.
\end{Proof}

\begin{defn}{Variation of Parameters}{varParamDef}
    Given a fundamental set $\{y_1, y_2\}$ for the homogeneous equation, we seek a particular solution of the form:
    \[
        y_p(t) = u_1(t)y_1(t) + u_2(t)y_2(t)
    \]
\end{defn}

\begin{thm}{Variation of Parameters Formula}{varParamThm}
    The particular solution is given by:
    \[
        y_p(t) = -y_1(t) \int \frac{y_2(t)g(t)}{W(t)} \, dt + y_2(t) \int \frac{y_1(t)g(t)}{W(t)} \, dt
    \]
    where $W(t) = W[y_1, y_2](t)$.
\end{thm}

\begin{Proof}
    We differentiate our guess $y_p = u_1 y_1 + u_2 y_2$:
    \[
        y_p' = u_1' y_1 + u_1 y_1' + u_2' y_2 + u_2 y_2'
    \]
    To simplify, we impose the condition: $u_1' y_1 + u_2' y_2 = 0$.
    Then $y_p' = u_1 y_1' + u_2 y_2'$. Differentiating again:
    \[
        y_p'' = u_1' y_1' + u_1 y_1'' + u_2' y_2' + u_2 y_2''
    \]
    Substitute $y_p, y_p', y_p''$ into the ODE. Since $y_1, y_2$ are homogeneous solutions, terms involving $u_1(\dots)$ and $u_2(\dots)$ vanish. We are left with:
    \[
        u_1' y_1' + u_2' y_2' = g(t)
    \]
    We now have a system of two linear equations for the unknowns $u_1', u_2'$:
    \[
        \begin{cases}
            u_1' y_1 + u_2' y_2 = 0 \\
            u_1' y_1' + u_2' y_2' = g(t)
        \end{cases}
    \]
    Using Cramer's Rule, the determinant is the Wronskian $W$. Thus:
    \[
        u_1' = \frac{\det\begin{pmatrix} 0 & y_2 \\ g(t) & y_2' \end{pmatrix}}{W} = \frac{-y_2 g(t)}{W}, \qquad u_2' = \frac{\det\begin{pmatrix} y_1 & 0 \\ y_1' & g(t) \end{pmatrix}}{W} = \frac{y_1 g(t)}{W}
    \]
    Integrating yields the result.
\end{Proof}

\begin{example}{Variation of Parameters}{varParamEx}
    Solve $y'' + y = \sec t$.

    \emph{Solution:}
    Homogeneous solution: $y_1 = \cos t, y_2 = \sin t$. $W = 1$.
    \[
        u_1(t) = -\int \sin t \sec t \, dt = -\int \tan t \, dt = -\ln|\sec t|
    \]
    \[
        u_2(t) = \int \cos t \sec t \, dt = \int 1 \, dt = t
    \]
    Particular solution: $y_p = (-\ln|\sec t|)\cos t + (t)\sin t$.
\end{example}

\chapter{Systems of Linear Equations}\label{cha:systems}

\subsection*{Summary}
While scalar equations describe single quantities, most physical models (interacting species, coupled oscillators, chemical reactions) involve multiple variables changing simultaneously. We model these as systems. We will focus on \emph{linear} systems, $\mathbf{x}' = A\mathbf{x}$, utilizing the powerful machinery of the Matrix Exponential $e^{At}$ to solve them elegantly, regardless of whether the matrix $A$ is diagonalizable or defective.

\section{Linear Systems and Geometry}

We define a first-order system of dimension $n$ using vector notation. Let $\mathbf{x}(t) \in \R^n$.

\begin{defn}{Linear System}{linearSystemDef}
    A first-order \emph{linear system} is an equation of the form:
    \[
        \mathbf{x}'(t) = A(t)\mathbf{x}(t) + \mathbf{f}(t)
    \]
    where $A(t)$ is an $n \times n$ matrix of continuous functions, and $\mathbf{f}(t)$ is a vector of continuous forcing functions. If $\mathbf{f}(t) \equiv \mathbf{0}$, the system is \emph{homogeneous}.
\end{defn}

Any $n$-th order scalar linear ODE $y^{(n)} + p_{n-1}y^{(n-1)} + \dots + p_0 y = 0$ can be converted into a first-order system of dimension $n$ by introducing variables $x_1 = y, x_2 = y', \dots, x_n = y^{(n-1)}$.

\begin{thm}{Existence and Uniqueness for Systems}{systemExistUnique}
    If $A(t)$ and $\mathbf{f}(t)$ are continuous on an interval $I$ containing $t_0$, then for any initial vector $\mathbf{x}_0 \in \R^n$, the IVP:
    \[
        \mathbf{x}' = A(t)\mathbf{x} + \mathbf{f}(t), \qquad \mathbf{x}(t_0) = \mathbf{x}_0
    \]
    has a unique solution $\mathbf{x}(t)$ defined on the entire interval $I$.
\end{thm}

\begin{Proof}
    We apply the Picard-Lindelöf Theorem for systems. We define the operator $T[\mathbf{x}](t) = \mathbf{x}_0 + \int_{t_0}^t (A(s)\mathbf{x}(s) + \mathbf{f}(s)) \, ds$.
    Since the matrix entries are continuous on a closed sub-interval, they are bounded by some $M$. The function $F(t, \mathbf{x}) = A(t)\mathbf{x} + \mathbf{f}(t)$ is Lipschitz continuous in $\mathbf{x}$ because:
    \[
        \| F(t, \mathbf{x}) - F(t, \mathbf{y}) \| = \| A(t)(\mathbf{x} - \mathbf{y}) \| \le \|A(t)\| \|\mathbf{x} - \mathbf{y}\| \le M \|\mathbf{x} - \mathbf{y}\|
    \]
    Thus, the contraction mapping principle guarantees a unique fixed point.
\end{Proof}

\begin{example}{Converting Scalar to System}{scalarToSystemEx}
    Convert the damped harmonic oscillator $y'' + 2y' + 2y = 0$ into a system.

    \emph{Solution:} Let $x_1 = y$ and $x_2 = y'$.
    Then $x_1' = x_2$.
    From the ODE, $x_2' = y'' = -2y' - 2y = -2x_2 - 2x_1$.
    The system is:
    \[
        \begin{pmatrix} x_1 \\ x_2 \end{pmatrix}' = \begin{pmatrix} 0 & 1 \\ -2 & -2 \end{pmatrix} \begin{pmatrix} x_1 \\ x_2 \end{pmatrix}
    \]
\end{example}

\section{The Matrix Exponential}

For the scalar case $x' = ax$, the solution is $x(t) = x_0 e^{at}$. We generalize this to systems $\mathbf{x}' = A\mathbf{x}$ (where $A$ is a constant matrix) by defining the exponential of a matrix.

\begin{defn}{Matrix Exponential}{matrixExpDef}
    For an $n \times n$ matrix $A$, the \emph{matrix exponential} is defined by the convergent power series:
    \[
        e^{A} = \sum_{k=0}^\infty \frac{A^k}{k!} = I + A + \frac{A^2}{2!} + \frac{A^3}{3!} + \dots
    \]
\end{defn}

This series converges absolutely for any matrix $A$ (under any matrix norm), making $e^A$ well-defined. Crucially, if $AB = BA$, then $e^{A+B} = e^A e^B$. However, if matrices do not commute, this identity fails.

\begin{thm}{Fundamental Matrix Solution}{fundMatrixSol}
    The matrix function $\Phi(t) = e^{At}$ is the unique solution to the matrix IVP:
    \[
        \Phi'(t) = A\Phi(t), \qquad \Phi(0) = I
    \]
    Consequently, the solution to the vector IVP $\mathbf{x}' = A\mathbf{x}, \mathbf{x}(0) = \mathbf{x}_0$ is given by:
    \[
        \mathbf{x}(t) = e^{At}\mathbf{x}_0
    \]
\end{thm}

\begin{Proof}
    We differentiate the power series term by term (justified by uniform convergence):
    \begin{align*}
        \frac{d}{dt} e^{At} &= \frac{d}{dt} \left( I + At + \frac{A^2 t^2}{2!} + \frac{A^3 t^3}{3!} + \dots \right) \\
        &= 0 + A + \frac{A^2 (2t)}{2!} + \frac{A^3 (3t^2)}{3!} + \dots \\
        &= A + A^2 t + \frac{A^3 t^2}{2!} + \dots \\
        &= A \left( I + At + \frac{A^2 t^2}{2!} + \dots \right) = A e^{At}
    \end{align*}
    At $t=0$, the series becomes $I + 0 + \dots = I$. By the Uniqueness Theorem, this is the unique solution.
\end{Proof}

\begin{example}{Computing Matrix Exponential (Diagonal Case)}{diagExpEx}
    Compute $e^{At}$ for $A = \begin{pmatrix} 2 & 0 \\ 0 & -3 \end{pmatrix}$.

    \emph{Solution:}
    Since $A$ is diagonal, $A^k = \begin{pmatrix} 2^k & 0 \\ 0 & (-3)^k \end{pmatrix}$.
    Substituting into the series:
    \[
        e^{At} = \sum \frac{t^k}{k!} \begin{pmatrix} 2^k & 0 \\ 0 & (-3)^k \end{pmatrix} = \begin{pmatrix} \sum \frac{(2t)^k}{k!} & 0 \\ 0 & \sum \frac{(-3t)^k}{k!} \end{pmatrix} = \begin{pmatrix} e^{2t} & 0 \\ 0 & e^{-3t} \end{pmatrix}
    \]
\end{example}

\section{Eigenvalues and Diagonalization}

In practice, we do not compute infinite sums. We use eigenvalues to diagonalize the matrix.

\begin{thm}{Diagonalizable Case}{diagThm}
    If $A$ has $n$ linearly independent eigenvectors $\mathbf{v}_1, \dots, \mathbf{v}_n$ with eigenvalues $\lambda_1, \dots, \lambda_n$, then $A$ is diagonalizable: $A = PDP^{-1}$, where $D = \operatorname{diag}(\lambda_1, \dots, \lambda_n)$ and $P = [\mathbf{v}_1 \dots \mathbf{v}_n]$.

    The solution is:
    \[
        e^{At} = P e^{Dt} P^{-1}
    \]
    Explicitly, the general solution is:
    \[
        \mathbf{x}(t) = c_1 e^{\lambda_1 t}\mathbf{v}_1 + \dots + c_n e^{\lambda_n t}\mathbf{v}_n
    \]
\end{thm}

\begin{Proof}
    Since $A = PDP^{-1}$, we have $A^k = (PDP^{-1})(PDP^{-1})\dots = PD^k P^{-1}$.
    Thus:
    \[
        e^{At} = \sum \frac{t^k}{k!} PD^k P^{-1} = P \left( \sum \frac{(Dt)^k}{k!} \right) P^{-1} = P e^{Dt} P^{-1}
    \]
    $e^{Dt}$ is simply the diagonal matrix with entries $e^{\lambda_i t}$.
\end{Proof}

\begin{example}{Real Distinct Eigenvalues}{realEigenEx}
    Solve $\mathbf{x}' = \begin{pmatrix} 1 & 1 \\ 4 & 1 \end{pmatrix}\mathbf{x}$.

    \emph{Solution:}
    Find eigenvalues: $\det(A - \lambda I) = (1-\lambda)^2 - 4 = \lambda^2 - 2\lambda - 3 = (\lambda - 3)(\lambda + 1)$.
    Roots: $\lambda_1 = 3, \lambda_2 = -1$.
    Eigenvectors:
    For $\lambda_1 = 3$: $\begin{pmatrix} -2 & 1 \\ 4 & -2 \end{pmatrix}\mathbf{v} = 0 \implies \mathbf{v}_1 = \begin{pmatrix} 1 \\ 2 \end{pmatrix}$.
    For $\lambda_2 = -1$: $\begin{pmatrix} 2 & 1 \\ 4 & 2 \end{pmatrix}\mathbf{v} = 0 \implies \mathbf{v}_2 = \begin{pmatrix} 1 \\ -2 \end{pmatrix}$.
    General Solution:
    \[
        \mathbf{x}(t) = c_1 e^{3t} \begin{pmatrix} 1 \\ 2 \end{pmatrix} + c_2 e^{-t} \begin{pmatrix} 1 \\ -2 \end{pmatrix}
    \]
\end{example}

\section{Defective Matrices and Jordan Forms}

What if $A$ does not have $n$ independent eigenvectors? This occurs when eigenvalues are repeated and the geometric multiplicity is less than the algebraic multiplicity.

We can decompose any matrix $A$ as $A = \lambda I + N$, where $N$ is nilpotent (meaning $N^k = 0$ for some $k$). Since $\lambda I$ and $N$ commute, $e^{At} = e^{\lambda I t}e^{Nt} = e^{\lambda t} e^{Nt}$. Because $N$ is nilpotent, the series for $e^{Nt}$ \emph{truncates} to a finite polynomial!

\begin{defn}{Generalized Eigenvector}{genEigenDef}
    A vector $\mathbf{v}$ is a \emph{generalized eigenvector} of rank $k$ corresponding to eigenvalue $\lambda$ if:
    \[
        (A - \lambda I)^k \mathbf{v} = \mathbf{0} \quad \text{but} \quad (A - \lambda I)^{k-1} \mathbf{v} \ne \mathbf{0}
    \]
\end{defn}

\begin{thm}{Solutions for Defective Matrices}{defectiveThm}
    If an eigenvalue $\lambda$ has algebraic multiplicity 2 but only 1 eigenvector $\mathbf{v}$, we find a second generalized eigenvector $\mathbf{w}$ such that:
    \[
        (A - \lambda I)\mathbf{w} = \mathbf{v}
    \]
    The two linearly independent solutions are:
    \[
        \mathbf{x}_1(t) = e^{\lambda t}\mathbf{v}
    \]
    \[
        \mathbf{x}_2(t) = e^{\lambda t}(\mathbf{w} + t(A-\lambda I)\mathbf{w}) = e^{\lambda t}(\mathbf{w} + t\mathbf{v})
    \]
\end{thm}

\begin{Proof}
    We compute $e^{At}\mathbf{w}$. Since $A\mathbf{w} = \lambda \mathbf{w} + \mathbf{v}$, we recall the decomposition $A = \lambda I + (A - \lambda I)$.
    \[
        e^{At}\mathbf{w} = e^{\lambda t} e^{(A - \lambda I)t} \mathbf{w} = e^{\lambda t} \left( I + t(A - \lambda I) + \frac{t^2}{2}(A-\lambda I)^2 + \dots \right) \mathbf{w}
    \]
    Since $\mathbf{w}$ is a generalized eigenvector of rank 2, $(A-\lambda I)^2 \mathbf{w} = (A-\lambda I)\mathbf{v} = 0$. The series terminates after the second term:
    \[
        e^{At}\mathbf{w} = e^{\lambda t}(\mathbf{w} + t\mathbf{v})
    \]
    This proves $\mathbf{x}_2$ is a solution.
\end{Proof}

\begin{example}{Repeated Eigenvalues}{repeatedEx}
    Solve $\mathbf{x}' = \begin{pmatrix} 1 & -1 \\ 1 & 3 \end{pmatrix}\mathbf{x}$.

    \emph{Solution:}
    Char. Poly: $(\lambda - 1)(\lambda - 3) + 1 = \lambda^2 - 4\lambda + 4 = (\lambda - 2)^2$. $\lambda = 2$.
    Eigenvector: $(A - 2I)\mathbf{v} = \begin{pmatrix} -1 & -1 \\ 1 & 1 \end{pmatrix}\mathbf{v} = 0 \implies \mathbf{v} = \begin{pmatrix} 1 \\ -1 \end{pmatrix}$.
    Only one eigenvector. We need a generalized eigenvector $\mathbf{w}$:
    \[
        (A - 2I)\mathbf{w} = \mathbf{v} \implies \begin{pmatrix} -1 & -1 \\ 1 & 1 \end{pmatrix} \begin{pmatrix} w_1 \\ w_2 \end{pmatrix} = \begin{pmatrix} 1 \\ -1 \end{pmatrix}
    \]
    $-w_1 - w_2 = 1$. Let $w_1 = 0$, then $w_2 = -1$. $\mathbf{w} = \begin{pmatrix} 0 \\ -1 \end{pmatrix}$.
    General Solution:
    \[
        \mathbf{x}(t) = c_1 e^{2t}\begin{pmatrix} 1 \\ -1 \end{pmatrix} + c_2 e^{2t} \left( \begin{pmatrix} 0 \\ -1 \end{pmatrix} + t \begin{pmatrix} 1 \\ -1 \end{pmatrix} \right)
    \]
\end{example}

\section{Phase Portraits and Stability}

For $2 \times 2$ systems, we can classify the qualitative behavior of solutions based on the trace $\tau = \operatorname{tr}(A)$ and determinant $\Delta = \det(A)$. The characteristic equation is $\lambda^2 - \tau\lambda + \Delta = 0$.

\begin{itemize}
    \item \emph{Saddle Point ($\Delta < 0$):} Real eigenvalues with opposite signs. Unstable.
    \item \emph{Node ($\Delta > 0, \tau^2 - 4\Delta > 0$):} Real eigenvalues with same sign.
    \begin{itemize}
        \item Stable (Sink) if $\tau < 0$.
        \item Unstable (Source) if $\tau > 0$.
    \end{itemize}
    \item \emph{Spiral ($\Delta > 0, \tau^2 - 4\Delta < 0$):} Complex eigenvalues $\alpha \pm i\beta$.
    \begin{itemize}
        \item Stable Spiral if $\tau < 0$ ($\alpha < 0$).
        \item Unstable Spiral if $\tau > 0$ ($\alpha > 0$).
        \item Center if $\tau = 0$ (Purely imaginary eigenvalues).
    \end{itemize}
\end{itemize}

\begin{example}{Classifying Stability}{stabilityEx}
    Classify the origin for the system with $A = \begin{pmatrix} 0 & 1 \\ -5 & -2 \end{pmatrix}$.

    \emph{Solution:}
    Trace $\tau = -2$. Determinant $\Delta = 5$.
    Discriminant: $\tau^2 - 4\Delta = 4 - 20 = -16 < 0$. Complex eigenvalues.
    Since $\tau = -2 < 0$, the real part is negative.
    Classification: \emph{Stable Spiral}.
\end{example}


\chapter{Qualitative Theory and Stability}\label{cha:stability}

\subsection*{Summary}
Most differential equations, particularly non-linear ones, cannot be solved in closed form. In this chapter, we abandon the search for explicit formulae and instead develop geometric and analytical tools to determine the qualitative behavior of solutions. We introduce the phase plane, define rigorous notions of stability (Lyapunov), and establish the relationship between a non-linear system and its linearization.

\section{Autonomous Systems and Phase Space}

We focus on \emph{autonomous} systems, where the governing laws do not change with time.

\begin{defn}{Autonomous System}{autonomousDef}
    A system of ODEs is \emph{autonomous} if it can be written as:
    \[
        \mathbf{x}' = \mathbf{f}(\mathbf{x})
    \]
    where $\mathbf{f}: U \to \R^n$ is a $C^1$ vector field on an open set $U \sse \R^n$. The independent variable $t$ does not appear explicitly on the right-hand side.
\end{defn}

\begin{defn}{Critical Point}{criticalPointDef}
    A point $\mathbf{x}_0$ is a \emph{critical point} (or equilibrium point) if $\mathbf{f}(\mathbf{x}_0) = \mathbf{0}$.

    If the system starts at a critical point ($\mathbf{x}(0) = \mathbf{x}_0$), the unique solution is the constant function $\mathbf{x}(t) \equiv \mathbf{x}_0$ for all $t$.
\end{defn}

The set of all solution curves (trajectories) in the domain $U$ forms the \emph{Phase Portrait}. For autonomous systems, trajectories never cross (by Uniqueness).

Intuitively, an equilibrium is ``stable'' if small disturbances do not grow over time. We formalize this using $\epsilon-\delta$ language.

\begin{defn}{Types of Stability}{stabilityTypesDef}
    Let $\mathbf{x}_0$ be a critical point.
    \begin{enumerate}
        \item \emph{Stable (in the sense of Lyapunov):} For every $\epsilon > 0$, there exists a $\delta > 0$ such that if $\|\mathbf{x}(0) - \mathbf{x}_0\| < \delta$, then $\|\mathbf{x}(t) - \mathbf{x}_0\| < \epsilon$ for all $t \ge 0$.
        \item \emph{Asymptotically Stable:} The point is stable, and there exists a $\delta_0 > 0$ such that if $\|\mathbf{x}(0) - \mathbf{x}_0\| < \delta_0$, then $\lim_{t \to \infty} \mathbf{x}(t) = \mathbf{x}_0$.
        \item \emph{Unstable:} The point is not stable.
    \end{enumerate}
\end{defn}

Stability means the solution stays close. Asymptotic stability means it stays close \emph{and} eventually converges to the equilibrium.


For a non-linear system $\mathbf{x}' = \mathbf{f}(\mathbf{x})$, the behavior near a critical point $\mathbf{x}_0$ is often well-approximated by the linear term of its Taylor series.
\[
    \mathbf{f}(\mathbf{x}) \approx \mathbf{f}(\mathbf{x}_0) + D\mathbf{f}(\mathbf{x}_0)(\mathbf{x} - \mathbf{x}_0) = J(\mathbf{x} - \mathbf{x}_0)
\]
where $J = D\mathbf{f}(\mathbf{x}_0)$ is the \emph{Jacobian matrix} evaluated at the critical point.

\begin{thm}{Principle of Linearized Stability}{linearizedStability}
    Let $\mathbf{x}' = \mathbf{f}(\mathbf{x})$ have a critical point at $\mathbf{0}$ (w.l.o.g.). Let $J = D\mathbf{f}(\mathbf{0})$.
    \begin{enumerate}
        \item If all eigenvalues of $J$ have strictly negative real parts ($\Re(\lambda_i) < 0$), then $\mathbf{0}$ is \emph{asymptotically stable}.
        \item If at least one eigenvalue has a strictly positive real part ($\Re(\lambda_i) > 0$), then $\mathbf{0}$ is \emph{unstable}.
        \item If some eigenvalues have zero real part (and none positive), the linearization is inconclusive (the center case).
    \end{enumerate}
\end{thm}

\begin{Proof}
    We sketch the proof for (1). Since all eigenvalues have negative real parts, there exists a specific inner product $\langle \cdot, \cdot \rangle_*$ and norm $\|\cdot\|_*$ such that $\langle J\mathbf{x}, \mathbf{x} \rangle_* \le -c \|\mathbf{x}\|_*^2$ for some $c > 0$.
    Consider the squared norm $V(t) = \|\mathbf{x}(t)\|_*^2$. Its derivative along trajectories is:
    \[
        \frac{d}{dt} \|\mathbf{x}\|_*^2 = 2\langle \mathbf{x}', \mathbf{x} \rangle_* = 2\langle J\mathbf{x} + \mathbf{g}(\mathbf{x}), \mathbf{x} \rangle_*
    \]
    where $\mathbf{g}(\mathbf{x})$ represents higher-order terms ($\|\mathbf{g}(\mathbf{x})\|/\|\mathbf{x}\| \to 0$).
    \[
        \frac{dV}{dt} \le -2c \|\mathbf{x}\|_*^2 + 2\|\mathbf{g}(\mathbf{x})\|_* \|\mathbf{x}\|_*
    \]
    Near $\mathbf{0}$, the quadratic decay dominated the higher-order terms, proving $V(t)$ decreases, so $\mathbf{x}(t) \to \mathbf{0}$.
\end{Proof}

\begin{example}{Linearization Analysis}{linearizationEx}
    Analyze the stability of the critical points for the pendulum equation with damping:
    \[
        \theta'' + \gamma \theta' + \sin \theta = 0, \quad \gamma > 0
    \]
    \emph{Solution:}
    Convert to a system: $x = \theta, y = \theta'$.
    \[
        \begin{cases} x' = y \\ y' = -\sin x - \gamma y \end{cases}
    \]
    Critical points occur where $y=0$ and $\sin x = 0$, so $x = n\pi$.
    The Jacobian is:
    \[
        J(x,y) = \begin{pmatrix} 0 & 1 \\ -\cos x & -\gamma \end{pmatrix}
    \]
    At $(0,0)$ (downward position): $J = \begin{pmatrix} 0 & 1 \\ -1 & -\gamma \end{pmatrix}$.
    Char. Poly: $\lambda^2 + \gamma \lambda + 1 = 0$. Roots $\lambda = \frac{-\gamma \pm \sqrt{\gamma^2 - 4}}{2}$. Real parts are negative. \emph{Asymptotically Stable}.

    At $(\pi, 0)$ (upward position): $J = \begin{pmatrix} 0 & 1 \\ 1 & -\gamma \end{pmatrix}$.
    Char. Poly: $\lambda^2 + \gamma \lambda - 1 = 0$. Roots $\lambda = \frac{-\gamma \pm \sqrt{\gamma^2 + 4}}{2}$. One root is positive. \emph{Unstable Saddle}.
\end{example}

\section{Lyapunov's Direct Method}

Linearization fails when eigenvalues have zero real parts. Lyapunov's Direct Method (or Second Method) allows us to determine stability without solving the equation, by generalizing the concept of energy.

\begin{defn}{Lyapunov Function}{lyapunovFuncDef}
    Let $\mathbf{x}^*$ be a critical point. A scalar function $V: U \to \R$ is a \emph{Lyapunov Function} if:
    \begin{enumerate}
        \item $V(\mathbf{x}^*) = 0$ and $V(\mathbf{x}) > 0$ for $\mathbf{x} \ne \mathbf{x}^*$ (Positive Definite).
        \item $\dot{V}(\mathbf{x}) = \nabla V \cdot \mathbf{f}(\mathbf{x}) \le 0$ for all $\mathbf{x} \in U \setminus \{\mathbf{x}^*\}$ (Orbital Derivative is Negative Semi-Definite).
    \end{enumerate}
    If the inequality in (2) is strict ($\dot{V} < 0$), it is a \emph{Strict Lyapunov Function}.
\end{defn}

\begin{thm}{Lyapunov Stability Theorem}{lyapunovThm}
    Let $\mathbf{x}^*$ be a critical point.
    \begin{enumerate}
        \item If there exists a Lyapunov Function $V$ in a neighborhood of $\mathbf{x}^*$, then $\mathbf{x}^*$ is \emph{Stable}.
        \item If there exists a Strict Lyapunov Function $V$, then $\mathbf{x}^*$ is \emph{Asymptotically Stable}.
    \end{enumerate}
\end{thm}

\begin{Proof}
    We prove (1). Let $\epsilon > 0$ be small enough that the ball $B_\epsilon(\mathbf{x}^*)$ is in the domain of $V$.
    Let $m = \min_{\|\mathbf{x}\| = \epsilon} V(\mathbf{x})$. Since $V$ is continuous and positive definite, $m > 0$.
    Define the set $U = \{\mathbf{x} \mid V(\mathbf{x}) < m\}$. Since $V(\mathbf{x}^*) = 0$ and $V$ is continuous, there exists a $\delta > 0$ such that $B_\delta(\mathbf{x}^*) \subset U$.

    Now, let a trajectory start at $\mathbf{x}_0$ with $\|\mathbf{x}_0 - \mathbf{x}^*\| < \delta$. Then $V(\mathbf{x}_0) < m$.
    Since $\dot{V} \le 0$, $V(\mathbf{x}(t))$ is non-increasing.
    Thus, $V(\mathbf{x}(t)) \le V(\mathbf{x}_0) < m$ for all $t \ge 0$.
    Since the boundary of $B_\epsilon$ has $V \ge m$, the trajectory can never reach the boundary $\|\mathbf{x}\| = \epsilon$.
    Therefore, $\|\mathbf{x}(t) - \mathbf{x}^*\| < \epsilon$ for all $t$. This is the definition of stability.
\end{Proof}

\begin{example}{Lyapunov Function}{lyapunovEx}
    Show that the origin is stable for the system:
    \[
        x' = -x - xy^2, \qquad y' = -y - yx^2
    \]
    \emph{Solution:}
    Try $V(x,y) = x^2 + y^2$. This is positive definite.
    Compute $\dot{V}$:
    \begin{align*}
        \dot{V} &= \frac{\partial V}{\partial x} x' + \frac{\partial V}{\partial y} y' \\
        &= 2x(-x - xy^2) + 2y(-y - yx^2) \\
        &= -2x^2 - 2x^2y^2 - 2y^2 - 2y^2x^2 \\
        &= -2(x^2 + y^2) - 4x^2y^2
    \end{align*}
    Clearly $\dot{V} < 0$ for all $(x,y) \ne (0,0)$. Thus, $V$ is a strict Lyapunov function, and the origin is asymptotically stable.
\end{example}

\section{Limit Cycles}

In non-linear systems, solutions can approach a closed periodic orbit called a \emph{limit cycle}.

\begin{thm}{Poincaré-Bendixson Theorem}{poincareBendixson}
    Suppose $R$ is a closed, bounded region in the plane ($\R^2$) that contains no critical points. If a trajectory $\mathbf{x}(t)$ is confined in $R$ for all $t \ge 0$, then $\mathbf{x}(t)$ must either be a periodic orbit or approach a periodic orbit as $t \to \infty$.
\end{thm}

\begin{Proof}
    (Sketch) The proof relies on the Jordan Curve Theorem (unique to $\R^2$). Since trajectories cannot cross, a bounded trajectory that doesn't terminate at a fixed point must eventually wind onto itself, forming a limit cycle.
\end{Proof}

\begin{example}{Existence of Limit Cycle}{limitCycleEx}
    Consider $r' = r(1-r^2), \theta' = 1$ (in polar coordinates).
    \begin{itemize}
        \item If $r < 1$, $r' > 0$ (radius increases).
        \item If $r > 1$, $r' < 0$ (radius decreases).
    \end{itemize}
    Consider the annulus $R = \{1/2 \le r \le 2\}$. Trajectories enter $R$ on both boundaries and never leave. Since there are no critical points (as $\theta' = 1 \ne 0$), any solution starting in $R$ must approach the limit cycle at $r=1$.
\end{example}

\chapter{Boundary Value Problems and Sturm-Liouville Theory}\label{cha:sturmLiouville}

\subsection*{Summary}
Up to this point, we have focused on Initial Value Problems (IVPs), where all conditions are specified at a single point $t_0$. In this chapter, we study \emph{Boundary Value Problems} (BVPs), where conditions are imposed at two different points. This shift leads to a fundamental difference: while linear IVPs always have unique solutions, BVPs is intimately tied to the spectrum of the differential operator. We develop the Sturm-Liouville theory, which generalizes the concept of diagonalization to differential operators, and introduce Green's functions for solving inhomogeneous problems.

\section{Introduction to Boundary Value Problems}

Consider the second-order linear equation on an interval $[a, b]$:
\[
    L[y] = y'' + p(x)y' + q(x)y = g(x), \qquad a < x < b
\]
subject to boundary conditions at the endpoints:
\[
    B_1[y] = \alpha_1 y(a) + \alpha_2 y'(a) = \gamma_1, \qquad B_2[y] = \beta_1 y(b) + \beta_2 y'(b) = \gamma_2
\]

\begin{defn}{Homogeneous BVP}{homogeneousBVPDef}
    A BVP is called \emph{homogeneous} if both the differential equation ($g(x) \equiv 0$) and the boundary conditions ($\gamma_1 = \gamma_2 = 0$) are homogeneous.
\end{defn}

The trivial solution $y(x) \equiv 0$ is always a solution to a homogeneous BVP. The central question of this chapter is: \emph{Under what conditions does a non-trivial solution exist?}

\begin{example}{Eigenvalues of a BVP}{bvpEigenEx}
    Consider the problem:
    \[
        y'' + \lambda y = 0, \qquad y(0) = 0, \quad y(\pi) = 0
    \]
    Depending on the parameter $\lambda$, the nature of the solution changes:
    \begin{enumerate}
        \item If $\lambda < 0$ (let $\lambda = -\mu^2$): $y(x) = c_1 \cosh \mu x + c_2 \sinh \mu x$. The boundary conditions force $c_1 = c_2 = 0$.
        \item If $\lambda = 0$: $y(x) = c_1 x + c_2$. Boundary conditions force $c_1 = c_2 = 0$.
        \item If $\lambda > 0$ (let $\lambda = \mu^2$): $y(x) = c_1 \cos \mu x + c_2 \sin \mu x$.
        $y(0)=0 \implies c_1 = 0$.
        $y(\pi) = c_2 \sin(\mu \pi) = 0$.
        For a non-trivial solution, we require $\sin(\mu \pi) = 0$, which implies $\mu = n$ for $n \in \mathbb{Z}^+$.
    \end{enumerate}
    Thus, non-trivial solutions (eigenfunctions) $y_n(x) = \sin(nx)$ exist only for the discrete spectrum $\lambda_n = n^2$.
\end{example}


To understand the existence of eigenvalues generally, we must define the ``transpose'' of a differential operator.

\begin{defn}{Formal Adjoint}{adjointDef}
    For the operator $L[y] = p_2(x)y'' + p_1(x)y' + p_0(x)y$, the \emph{formal adjoint} $L^*$ is defined by:
    \[
        L^*[u] = \frac{d^2}{dx^2}(p_2 u) - \frac{d}{dx}(p_1 u) + p_0 u
    \]
    The operator $L$ is \emph{formally self-adjoint} if $L = L^*$.
\end{defn}

A straightforward calculation shows that $L$ is self-adjoint if and only if $p_1(x) = p_2'(x)$. This motivates the \emph{Sturm-Liouville form}.

\begin{defn}{Sturm-Liouville Form}{slFormDef}
    A second-order differential operator $L$ is in \emph{Sturm-Liouville form} if:
    \[
        L[y] = -\frac{d}{dx}\left[ p(x)\frac{dy}{dx} \right] + q(x)y
    \]
    where $p(x) > 0$. This operator is formally self-adjoint.
\end{defn}

\emph{Note:} Any linear operator $P(x)y'' + Q(x)y' + R(x)y$ can be multiplied by the integrating factor $\mu(x) = \frac{1}{P(x)}e^{\int Q/P dx}$ to be converted into Sturm-Liouville form.

\begin{thm}{Lagrange's Identity}{lagrangeIdentity}
    For a self-adjoint operator $L$ in Sturm-Liouville form:
    \[
        u L[v] - v L[u] = \frac{d}{dx} \left[ p(x)(u v' - u' v) \right]
    \]
    Integrating over $[a,b]$ yields Green's Formula:
    \[
        \langle L[u], v \rangle - \langle u, L[v] \rangle = \left[ p(x)(u v' - u' v) \right]_a^b
    \]
\end{thm}

If $u$ and $v$ satisfy \emph{separated boundary conditions} (e.g., $\alpha_1 u(a) + \alpha_2 u'(a) = 0$), the boundary term vanishes, and we obtain the crucial symmetry property: $\langle L[u], v \rangle = \langle u, L[v] \rangle$.

\section{The Sturm-Liouville Theorems}

We now consider the eigenvalue problem:
\[
    L[y] = \lambda w(x) y
\]
subject to homogeneous separated boundary conditions.

\begin{thm}{Sturm-Liouville Theorem}{slTheorem}
    For a regular Sturm-Liouville problem:
    \begin{enumerate}
        \item All eigenvalues $\lambda$ are real.
        \item Eigenfunctions $\phi_n, \phi_m$ corresponding to distinct eigenvalues are \emph{orthogonal} with respect to the weight $w(x)$:
        \[
            \int_a^b \phi_n(x) \phi_m(x) w(x) \, dx = 0 \quad \text{if } n \ne m
        \]
        \item The eigenvalues form an infinite sequence $\lambda_1 < \lambda_2 < \dots$ with $\lim_{n\to\infty} \lambda_n = \infty$.
        \item The eigenfunctions form a \emph{complete basis} for the space of piecewise continuous functions $L^2_w[a,b]$.
    \end{enumerate}
\end{thm}

\begin{Proof}
    \emph{Real Eigenvalues:} Let $L[u] = \lambda w u$. Take the complex conjugate $L[\bar{u}] = \bar{\lambda} w \bar{u}$.
    Using Green's Formula (with boundary terms vanishing):
    \[
        (\lambda - \bar{\lambda}) \langle w u, \bar{u} \rangle = \langle L[u], \bar{u} \rangle - \langle u, L[\bar{u}] \rangle = 0
    \]
    Since $w(x) > 0$, $\langle w u, \bar{u} \rangle = \int w |u|^2 > 0$, implying $\lambda = \bar{\lambda}$.

    \emph{Orthogonality:} Let $L[\phi_n] = \lambda_n w \phi_n$ and $L[\phi_m] = \lambda_m w \phi_m$.
    \[
        (\lambda_n - \lambda_m) \langle w \phi_n, \phi_m \rangle = \langle L[\phi_n], \phi_m \rangle - \langle \phi_n, L[\phi_m] \rangle = 0
    \]
    If $\lambda_n \ne \lambda_m$, the inner product must vanish.
\end{Proof}

\begin{thm}{Eigenfunction Expansion}{completenessThm}
    Any piecewise smooth function $f(x)$ on $[a,b]$ can be represented as:
    \[
        f(x) \sim \sum_{n=1}^\infty c_n \phi_n(x)
    \]
    where the Fourier coefficients are given by the projection:
    \[
        c_n = \frac{\langle f, \phi_n \rangle_w}{\langle \phi_n, \phi_n \rangle_w} = \frac{\int_a^b f(x)\phi_n(x)w(x) \, dx}{\int_a^b \phi_n(x)^2 w(x) \, dx}
    \]
\end{thm}

\section{The Fredholm Alternative}

We now address the inhomogeneous equation $L[y] - \lambda w(x) y = f(x)$. The solvability depends on whether $\lambda$ is an eigenvalue.

\begin{thm}{Fredholm Alternative}{fredholmAlt}
    Consider the inhomogeneous BVP $L[y] - \sigma w(x)y = f(x)$.
    \begin{enumerate}
        \item \textbf{Non-Resonant Case:} If $\sigma$ is \emph{not} an eigenvalue of the homogeneous problem, there exists a unique solution for any $f(x)$.
        \item \textbf{Resonant Case:} If $\sigma = \lambda_k$ is an eigenvalue with eigenfunction $\phi_k$, a solution exists \emph{if and only if} $f$ is orthogonal to the eigenfunction:
        \[
            \int_a^b f(x)\phi_k(x) w(x) \, dx = 0
        \]
        If this condition holds, there are infinitely many solutions (of the form $y_p + c\phi_k$).
    \end{enumerate}
\end{thm}

\section{Green's Functions}

In the non-resonant case, the unique solution can be represented by an integral kernel called the Green's function.

\begin{defn}{Green's Function}{greensFuncDef}
    The \emph{Green's Function} $G(x, \xi)$ for the operator $L$ is the solution to the distributional equation:
    \[
        L[G(x, \xi)] = \delta(x - \xi)
    \]
    subject to the homogeneous boundary conditions.
\end{defn}

To construct $G(x, \xi)$ explicitly, we solve the homogeneous equation $L[y]=0$ on the intervals $[a, \xi)$ and $(\xi, b]$.
Let $y_1(x)$ satisfy the boundary condition at $a$, and $y_2(x)$ satisfy the condition at $b$. We assume:
\[
    G(x, \xi) = \begin{cases}
        c_1(\xi) y_1(x) & a \le x < \xi \\
        c_2(\xi) y_2(x) & \xi < x \le b
    \end{cases}
\]
We require two conditions at $x = \xi$:
\begin{enumerate}
	\item   \emph{Continuity:} $c_1 y_1(\xi) = c_2 y_2(\xi)$.
	\item  \emph{Jump in Derivative:} Integrating $-\frac{d}{dx}[p y'] = \delta$ from $\xi-\epsilon$ to $\xi+\epsilon$ implies a jump of $-1/p(\xi)$:
    \[
        \frac{\partial G}{\partial x}(\xi^+, \xi) - \frac{\partial G}{\partial x}(\xi^-, \xi) = -\frac{1}{p(\xi)}
    \]
\end{enumerate}

Solving this system for $c_1, c_2$ using the Wronskian $W(\xi) = y_1 y_2' - y_2 y_1'$ yields the symmetric form:

\begin{thm}{Structure of the Green's Function}{greensStruct}
    The Green's function is given by:
    \[
        G(x, \xi) = \frac{1}{p(\xi)W(\xi)} \begin{cases}
            y_1(x)y_2(\xi) & a \le x \le \xi \\
            y_1(\xi)y_2(x) & \xi \le x \le b
        \end{cases}
    \]
    Note that $p(\xi)W(\xi)$ is a constant. The solution to $L[y] = f(x)$ is then:
    \[
        y(x) = \int_a^b G(x, \xi)f(\xi) \, d\xi
    \]
\end{thm}



\chapter{Series Solutions and Special Functions}\label{cha:seriesSolutions}

\subsection*{Summary}
Throughout the previous chapters, we primarily dealt with linear equations with \emph{constant} coefficients. However, the Partial Differential Equations of mathematical physics (Heat, Wave, Schrödinger) often reduce to ODEs with \emph{variable} coefficients when solved in non-Cartesian coordinates (cylindrical, spherical). These equations rarely have elementary solutions ($e^x, \sin x$). In this chapter, we develop the method of power series (Frobenius Method) to define and compute these new ``Special Functions,'' specifically Bessel functions and Legendre polynomials.

\section{Analytic Functions and Ordinary Points}

We consider the homogeneous linear second-order equation:
\begin{equation}\label{eq:varCoeff}
    P(x)y'' + Q(x)y' + R(x)y = 0
\end{equation}
Dividing by $P(x)$, we obtain the standard form $y'' + p(x)y' + q(x)y = 0$. The behavior of solutions depends entirely on the behavior of $p(x)$ and $q(x)$ at a given point $x_0$.

\begin{defn}{Ordinary and Singular Points}{ordinarySingularDef}
    Let $p(x) = Q(x)/P(x)$ and $q(x) = R(x)/P(x)$. A point $x_0$ is called an \emph{ordinary point} if both $p(x)$ and $q(x)$ are analytic at $x_0$ (i.e., they have convergent Taylor series expansions in a neighborhood of $x_0$).

	\qn
    If either function is not analytic at $x_0$, the point is called a \emph{singular point}.
\end{defn}

Since polynomials are analytic everywhere, if $P, Q, R$ are polynomials, singular points only occur at the roots of $P(x) = 0$.

\begin{thm}{Existence of Power Series Solutions}{powerSeriesThm}
    If $x_0$ is an ordinary point of the differential equation, then there exist two linearly independent solutions of the form:
    \[
        y(x) = \sum_{n=0}^\infty a_n (x - x_0)^n
    \]
    The radius of convergence of this series is at least as large as the distance from $x_0$ to the nearest singular point (in the complex plane).
\end{thm}

\begin{Proof}
    (Sketch) Since $p(x)$ and $q(x)$ are analytic, they can be represented by power series. We assume $y(x) = \sum a_n (x-x_0)^n$, differentiate term-by-term, and substitute into the ODE. Equating coefficients of like powers of $(x-x_0)$ yields a recurrence relation determining $a_n$ for $n \ge 2$ in terms of the arbitrary constants $a_0, a_1$. The convergence is guaranteed by the theory of analytic functions in complex analysis.
\end{Proof}

\begin{example}{Legendre's Equation}{legendreEx}
    Consider Legendre's equation, which arises in spherical coordinates:
    \[
        (1-x^2)y'' - 2xy' + \alpha(\alpha+1)y = 0
    \]
    Here $P(x) = 1-x^2$. The singular points are $x = \pm 1$. The point $x_0 = 0$ is an ordinary point.

    \emph{Solution:} Let $y = \sum a_n x^n$.
    \[
        y' = \sum n a_n x^{n-1}, \qquad y'' = \sum n(n-1)a_n x^{n-2}
    \]
    Substituting into the ODE:
    \[
        (1-x^2)\sum_{n=2}^\infty n(n-1)a_n x^{n-2} - 2x \sum_{n=1}^\infty n a_n x^{n-1} + \alpha(\alpha+1)\sum_{n=0}^\infty a_n x^n = 0
    \]
    Shifting indices to align powers of $x^k$:
    \[
        \sum_{k=0}^\infty \left[ (k+2)(k+1)a_{k+2} - (k(k-1) + 2k - \alpha(\alpha+1))a_k \right] x^k = 0
    \]
    The recurrence relation is:
    \[
        a_{k+2} = -\frac{(\alpha - k)(\alpha + k + 1)}{(k+2)(k+1)} a_k
    \]
    Notice that if $\alpha = n$ is a non-negative integer, the numerator vanishes when $k=n$, causing the series to terminate. These terminating solutions are the \emph{Legendre Polynomials} $P_n(x)$.
\end{example}

\section{Regular Singular Points and the Method of Frobenius}

Many physical potentials (like the Coulomb potential $1/r$) introduce singularities at the origin. Standard power series fail here. We need a more robust method.

\begin{defn}{Regular Singular Point}{regularSingularDef}
    A singular point $x_0$ is called a \emph{regular singular point} if the singularities are no worse than a pole of order 1 for $p(x)$ and order 2 for $q(x)$. Specifically, if:
    \[
        \lim_{x \to x_0} (x-x_0)p(x) \quad \text{and} \quad \lim_{x \to x_0} (x-x_0)^2 q(x)
    \]
    are both finite (analytic), then $x_0$ is regular.
\end{defn}

For such points, we assume a solution of the form:
\[
    y(x) = (x-x_0)^r \sum_{n=0}^\infty a_n (x-x_0)^n = \sum_{n=0}^\infty a_n (x-x_0)^{n+r}, \quad a_0 \ne 0
\]
This is called the \emph{Frobenius Method}. The exponent $r$ is determined by the ``Indicial Equation.''

\begin{thm}{Frobenius Theorem}{frobeniusThm}
    Let $x=0$ be a regular singular point. Let $p_0 = \lim_{x \to 0} xp(x)$ and $q_0 = \lim_{x \to 0} x^2q(x)$. The \emph{Indicial Equation} is:
    \[
        r(r-1) + p_0 r + q_0 = 0
    \]
    Let $r_1, r_2$ be the roots with $\Re(r_1) \ge \Re(r_2)$.
    \begin{enumerate}
        \item There always exists one solution $y_1(x) = |x|^{r_1} \sum a_n x^n$.
        \item If $r_1 - r_2$ is not an integer, the second solution is $y_2(x) = |x|^{r_2} \sum b_n x^n$.
        \item If $r_1 - r_2$ is an integer (or zero), the second solution may involve a logarithm:
        \[
            y_2(x) = C y_1(x) \ln|x| + |x|^{r_2} \sum b_n x^n
        \]
    \end{enumerate}
\end{thm}


The most important application of the Frobenius method is solving Bessel's equation, which arises when separating variables in cylindrical coordinates (Laplacian in 2D or 3D).

\begin{example}{Bessel's Equation of Order $\nu$}{besselEx}
    \[
        x^2 y'' + x y' + (x^2 - \nu^2)y = 0, \qquad \nu \ge 0
    \]
    Here $p(x) = 1/x$ and $q(x) = (x^2-\nu^2)/x^2$. $x=0$ is a regular singular point.

    \emph{Indicial Equation:}
    Multiply by $x^2$ to identify the leading terms: $x^2 y'' + xy' - \nu^2 y \approx 0$ near $x=0$. This is Cauchy-Euler!
    The characteristic equation is $r(r-1) + r - \nu^2 = 0 \implies r^2 - \nu^2 = 0 \implies r = \pm \nu$.

    \emph{Solution Construction:}
    Using $r_1 = \nu$, the Frobenius method yields the \emph{Bessel Function of the First Kind}, denoted $J_\nu(x)$:
    \[
        J_\nu(x) = \sum_{n=0}^\infty \frac{(-1)^n}{n! \Gamma(n+\nu+1)} \left(\frac{x}{2}\right)^{2n+\nu}
    \]
    For the second solution, since $r_1 - r_2 = 2\nu$ might be an integer, we encounter the logarithmic case. The second linearly independent solution is the \emph{Bessel Function of the Second Kind}, $Y_\nu(x)$, which blows up at $x=0$.
\end{example}

\section*{Sturm-Liouville Properties of Special Functions}

Why do we care about these complicated series? Because they form orthogonal bases, just like sines and cosines.

Recall Bessel's equation: $(x y')' + (x - \frac{\nu^2}{x})y = 0$. This is almost Sturm-Liouville form with $p(x)=x, q(x)=-\nu^2/x, \lambda=1, w(x)=x$.
Consider the BVP on $[0, R]$ with $y(R)=0$. The solutions are $J_\nu(\lambda_k x)$ where $\lambda_k$ are roots of $J_\nu(\lambda R) = 0$.

\begin{thm}{Orthogonality of Bessel Functions}{besselOrthog}
    Let $\{\lambda_k\}$ be the positive zeros of $J_\nu(x)$. Then the functions $\phi_k(x) = J_\nu(\lambda_k x/R)$ are orthogonal on $[0, R]$ with respect to the weight $w(x) = x$:
    \[
        \int_0^R J_\nu\left(\frac{\lambda_k x}{R}\right) J_\nu\left(\frac{\lambda_m x}{R}\right) x \, dx = 0 \qquad \text{if } k \ne m
    \]
\end{thm}

This orthogonality allows us to expand any function $f(r)$ on a disk into a Fourier-Bessel series:
\[
    f(r) = \sum_{k=1}^\infty c_k J_\nu\left(\frac{\lambda_k r}{R}\right)
\]
This is the fundamental mechanism for solving heat flow on a circular plate or vibrations of a drum head.



\end{document}

